\begin{textsample}{POS Dim 2 – 2017\_c6 – Score 1.00 – 2017\_c6/enade\_2017\_ci\_comp\_7c2b2a3864.txt}  \label{ex:f2_pos_269}
CAPÍTULO 6 ANÁLISE TÉCNICA DA PROVA Este capítulo tem por objetivo apresentar o desempenho dos estudantes concluintes de Ciência da Computação ( Licenciatura/Bacharelado ) no Enade/2017.
Para isso, foram calculadas as estatísticas básicas da prova em seu todo ( seção 6.1.1 ), bem como as estatísticas relacionadas aos Componentes de Formação Geral ( seção 6.1.2 ) e de Conhecimento Específico da Área ( seção 6.1.3 ).
Dadas as suas características, foram analisadas, em separado, as questões objetivas ( seção 6.2 ) e as questões discursivas ( seção 6.3 ).
Para as questões objetivas, foram disponibilizados os índices de facilidade e de discriminação ponto-bisserial, também em separado, para os Componentes de Formação Geral ( seções 6.2.1 e 6.2.2 ) e de Conhecimento Específico ( seções 6.2.3 e 6.2.4 ).
De cada componente, uma das questões foi escolhida para exemplificar a análise gráfica, relacionando as alternativas escolhidas pelos estudantes ( inclusive o gabarito ) com o número de acertos no componente.
O Anexo I apresenta a íntegra da análise gráfica para todas as questões objetivas.
Para cada uma das questões discursivas, os conteúdos dos tipos mais comuns de respostas dos estudantes são apresentados e comparados com o padrão de respostas esperado ( ver Anexo VIII com o padrão de respostas ).
Tomando como base as duas questões discursivas do Componente de Formação Geral, a seção 6.3.1.6 apresenta comentários sobre a correção das respostas com respeito à Língua Portuguesa.
As tabelas apresentam as seguintes estatísticas das notas20 : média do desempenho na prova, erro padrão da média, desvio padrão, nota mínima, mediana e nota máxima para cada um de seus componentes.
Tais estatísticas contemplam o total de estudantes concluintes da Área de Ciência da Computação ( Licenciatura/Bacharelado ) inscritos e presentes à prova do Enade/2017, tendo em vista agregações, ou por Grandes Regiões e o país como um todo, ou por Categoria Administrativa, Organização Acadêmica da IES e Modalidade de Ensino.
Em relação aos gráficos de distribuição de notas, o intervalo considerado foi de 10 unidades, aberto à esquerda e fechado à direita, com exceção do primeiro intervalo, [ 0 ; 10 ], fechado em ambos os extremos.
Para os gráficos de distribuição das notas das questões discursivas, foram consideradas mais duas categorias : questão em branco21 e nota zero. 20 Uma definição dessas estatísticas pode ser encontrada no Glossário. 21 Nesse grupo estão incluídas também as respostas classificadas como nulas ou desconsideradas. 6.1 ESTATÍSTICAS BÁSICAS DA PROVA Esta seção apresenta estatísticas selecionadas e histogramas da nota geral ( 6.1.1 ) e de cada componente : Formação Geral ( 6.1.2 ) e Conhecimento Específico ( 6.1.3 ).
São também apresentadas estatísticas selecionadas de subpopulações caracterizadas por Grande Região, Categoria Administrativa, Organização Acadêmica e Modalidade de Ensino. 6.1.1 Estatísticas Básicas Gerais A Tabela 6.1 apresenta as Estatísticas Básicas da prova, por Grande Região segundo Modalidade de Ensino e Habilitação, dos estudantes concluintes de Ciência da Computação ( Licenciatura/Bacharelado ).
As análises serão feitas, primeiramente, para os cursos de Licenciatura e, em seguida, para os de Bacharelado.
Cada habilitação é desagregada internamente em cursos Presenciais e cursos a Distância ( quando pertinente ).
A Média das notas da prova, como um todo, dos alunos de Licenciatura ( nas seções seguintes serão analisados os Componentes de Formação Geral e de Conhecimento Específico ), foi 40,8, sendo que os alunos da região Centro-Oeste obtiveram a Média mais baixa ( 36,0 ), e os da região Sul obtiveram a Média mais alta ( 47,8 ).
As demais médias foram : 37,7 na região Norte, 42,2 na região Nordeste e 42,0 na região Sudeste.
Já a Média das notas para os alunos de cursos Presenciais de Licenciatura foi um pouco mais alta, 41,1, sendo que os alunos da região Centro-Oeste obtiveram a Média mais baixa ( 37,0 ), e os da região Sul, a Média mais alta ( 47,8 ).
As demais médias foram : 43,1 na região Nordeste, 42,5 na região Sudeste e 37,9 na região Norte.
O Desvio padrão para o Brasil, como um todo, foi 13,1, sendo o maior Desvio padrão encontrado na região Centro-Oeste ( 15,4 ) e o menor, na região Norte ( 10,8 ), indicando uma dispersão um pouco menor das notas dessa última região.
Para as Licenciaturas Presenciais, a região que obteve a maior nota Máxima foi a Sul ( 79,6 ), ao passo que a região que atingiu a menor nota Máxima foi a Norte ( 69,5 ).
A Mediana do Brasil, como um todo, foi 41,6, sendo a maior Mediana obtida na região Sul ( 47,9 ), e a menor, na região Centro-Oeste ( 33,7 ).
A nota Mínima foi 6,5 na região Nordeste, 8,6 na Sudeste, 13,1 no Norte, 13,8 no Centro-Oeste e 20,8 no Sul.
Considerando -se as notas, segundo Grande Região, observa -se que existe diferença estatisticamente significativa ao nível de 95 \% entre a Média da região Centro-Oeste ( 37,0 ), a menor das Licenciaturas Presenciais dentre as Grandes Regiões, e as duas maiores médias, obtidas nas regiões Nordeste ( 43,1 ) e Sul (47,8)22.
Já a Média das notas da prova dos alunos de cursos a Distância de Licenciatura, foi 40,1, mais baixa do que os alunos de cursos Presenciais, sendo que os alunos da região Centro-Oeste obtiveram a Média mais baixa ( 31,9 ), e os da região Sul obtiveram a Média mais alta ( 51,6 ).
Na região Norte, a Média foi 32,6, na região Nordeste, 40,6, e na região Sudeste, 41,8.
O Desvio padrão para o Brasil como um todo foi 13,6, sendo o maior Desvio padrão encontrado na região Sudeste ( 14,9 ) e o menor, na região Sul ( 0,6 ).
A região que obteve a maior nota Máxima foi a Sudeste ( 75,2 ), ao passo que a região que atingiu a menor nota Máxima foi a Norte ( 43,8 ).
A Mediana do Brasil como um todo foi 38,8, sendo a maior Mediana obtida na região Sul ( 51,6 ), e a menor obtida na região Centro-Oeste ( 32,2 ).
A nota Mínima foi zero, obtida na região Nordeste.
Nas demais regiões as notas mínimas foram : 7,7 na região Sudeste, 8,3 na região Centro-Oeste, 18,1 na região Norte e 51,1 na região Sul.
Considerando -se as notas segundo Grande Região, observa -se que existe diferença estatisticamente significativa ao nível de 95 \% entre a menor Média, obtida na região Centro-Oeste ( 31,9 ), e as duas maiores médias, obtidas nas regiões Sul ( 51,6 ) e Sudeste ( 41,8 ).
Cumpre notar que ainda que as notas dos cursos presenciais sejam comparáveis com as dos cursos a Distância dentro de uma mesma Habilitação, não é possível comparar as notas obtidas pelos alunos de Licenciatura e Bacharelado, já que as provas são, em parte, diferentes.
A Média das notas da prova- como um todo- para os alunos de cursos de Bacharelado, foi 43,3, sendo que os alunos da região Norte obtiveram a Média mais baixa ( 38,4 ), e os da região Sul- a Média mais alta ( 46,2 ).
As demais médias foram : 44,0 na região Nordeste, 42,6 na região Sudeste e 41,3 na região Centro-Oeste.
Cabe observar que não houve participação de estudantes de Bacharelados a Distância na área de Ciência da Computação.
Assim, as estatísticas apresentadas para o Total dos Bacharelados são as mesmas que constam para os Bacharelados Presenciais.
A região que obteve a maior nota Máxima foi a Sudeste ( 95,1 ) e a que obteve a menor nota Máxima foi a Norte ( 76,6 ).
A Mediana do Brasil, como um todo, foi 42,3, sendo a maior Mediana obtida na região Sul ( 45,7 ), e a menor obtida na Norte ( 37,0 ).
A nota Mínima só não foi zero na região Norte, na qual essa nota foi 6,1.
Considerando -se as notas, segundo Grande Região, observa -se que existe diferença estatisticamente significativa ao nível de 95 \% entre a menor Média, obtida na região Norte ( 38,4 ), e as médias de todas as demais regiões.
O comportamento das notas dos estudantes de todo o Brasil pode ser observado no Gráfico 6.1, que apresenta um histograma com a distribuição das notas, segundo a habilitação.
Essa é uma distribuição unimodal para ambas as Habilitações, com as modas nos intervalos ( 40 ; 50 ] e ( 30 ; 40 ], respectivamente, para os alunos de cursos de Licenciatura e de Bacharelado.
A Tabela 6.2 apresenta informações referentes às Estatísticas Básicas da prova, desagregadas por Categoria Administrativa e por Organização Acadêmica, segundo a Habilitação.
Os estudantes de Licenciatura de IES Públicas obtiveram Média igual a Média nacional para esta Habilitação ( 40,8 ), e os das IES Privadas obtiveram Média um pouco mais baixa ( 40,7 ) que a Média nacional para essa Habilitação ( 40,8 ).
Observa -se que não existe diferença estatisticamente significativa entre as médias das notas das IES Públicas e Privadas.
O Desvio padrão para as IES Públicas ( 13,0 ) foi inferior ao do Brasil, como um todo ( 13,3 ), indicando uma dispersão menor das notas nessa Categoria Administrativa do que para o conjunto de alunos de cursos de Licenciatura.
No tocante à Organização Acadêmica, alunos de Licenciatura de Centros Universitários obtiveram nota Média mais alta ( 45,7 ) do que os de Universidades ( 41,3 ), que, por sua vez, obtiveram nota Média mais alta do que os de CEFET/IFET ( 39,2 ) e de Faculdades ( 33,8 ), com diferenças estatisticamente significativas ao nível de 95 \% entre essa última Média, a mais baixa, e todas as demais.
Também para o Bacharelado, em relação à Categoria Administrativa, os estudantes das IES Públicas obtiveram Média mais alta ( 48,4 ), e os das IES Privadas obtiveram Média mais baixa ( 38,8 ) que a Média nacional para essa habilitação ( 43,3 ).
Nesse caso, observa -se que existe diferença estatisticamente significativa entre as médias das notas das IES Públicas e Privadas.
O Desvio padrão para as IES Públicas ( 13,8 ) e para IES Privadas ( 11,8 ) foi, respectivamente, superior e inferior ao do Brasil, como um todo ( 13,6 ).
No tocante à Organização Acadêmica, alunos de Universidades e dos CEFET/IFET obtiveram média mais alta ( 45,1 ) do que os de Centros Universitários ( 39,6 ) e os de Faculdades ( 35,2 ).
Só não há diferença estatisticamente significativa ao nível de 95 \% entre as Médias das Universidades e as dos CEFET/IFET. 6.1.2 Estatísticas Básicas no Componente de Formação Geral A Tabela 6.3 apresenta as Estatísticas Básicas em relação ao componente da prova que avalia a Formação Geral dos estudantes concluintes, por Grande Região, segundo Habilitação e Modalidade de ensino.
Os alunos de cursos de Licenciatura obtiveram desempenho médio de 49,6.
Quanto à variabilidade, o Desvio padrão das notas dos estudantes do Brasil, como um todo, foi 17,5.
A maior Média foi obtida na região Sul ( 56,2 ), e a menor, na região Centro-Oeste ( 46,0 ).
As demais médias foram : 47,7 na região Norte e 50,4 nas regiões Nordeste e Sudeste.
Já os alunos de cursos Presenciais de Licenciatura obtiveram desempenho um pouco melhor ( 49,7 ) do que o agregado das duas modalidades.
Quanto à variabilidade, o Desvio padrão das notas dos estudantes das Licenciaturas Presenciais do Brasil, como um todo, foi 16,8.
A maior Média foi obtida na região Sul ( 55,9 ), e a menor, na região Centro-Oeste ( 47,2 ).
As demais Médias foram : 47,5 na região Norte, 48,6 na região Sudeste e 51,1 na região Nordeste.
Já o maior Desvio padrão foi obtido na região Centro-Oeste ( 18,2 ), e o menor, na região Sul ( 14,6 ).
Os demais desvios padrões foram : 16,1 na região Nordeste, 16,3 na região Norte e 17,9 na região Sudeste.
A maior nota no Componente de Formação Geral da prova de alunos de cursos presenciais de Licenciatura foi 96,0, obtida por, pelo menos, um aluno na região Norte.
A menor nota Máxima foi obtida na região Sul ( 82,9 ).
A Mediana do Brasil, como um todo, foi 50,4, sendo a menor Mediana encontrada na região Centro-Oeste ( 46,7 ), e a maior encontrada na região Sul ( 56,8 ).
A nota Mínima dos estudantes das Licenciaturas Presenciais nessa parte foi zero no Sudeste.
Na região Norte, a nota Mínima foi 6,6, nas regiões Nordeste e Centro-Oeste a nota Mínima foi 7,5 e na região Sul, 17,3.
Considerando -se as notas, segundo Grande Região, observa -se só não há diferença estatisticamente significativa entre a maior Média das notas do Componente de Formação Geral, obtidas na região Sul ( 55,9 ) e a nota Média da região Nordeste ( 51,1 ).
Já a Média das notas do Componente de Formação Geral dos alunos de cursos de Licenciatura a Distância, foi 49,5, mais baixa do que a dos alunos de cursos presenciais, sendo que os alunos da região Centro-Oeste obtiveram a Média mais baixa ( 41,2 ), e os da região Sul obtiveram a Média mais alta ( 70,3 ).
As demais médias foram : 51,9 na região Norte, 49,1 na região Nordeste e 51,3 na região Sudeste.
O Desvio padrão para o Brasil como um todo foi 19,0, sendo o maior Desvio padrão encontrado na região Sul ( 21,8 ) e o menor, na região Norte ( 11,4 ).
A região que obteve a maior nota Máxima foi a Nordeste ( 96,4 ), ao passo que a região que atingiu a menor nota Máxima foi a Norte ( 66,6 ).
A Mediana do Brasil como um todo foi 49,8, sendo a maior Mediana obtida na região Sul ( 70,3 ), e a menor obtida na região Centro-Oeste ( 44,1 ).
A nota Mínima foi zero nas regiões Nordeste e Centro-Oeste.
Na região Norte a nota Mínima foi 36,8, na Sudeste foi 7,5 e na região Sul, 54,8.
Considerando -se as notas segundo Grande Região, observa -se que só existe diferença estatisticamente significativa ao nível de 95 \% entre a menor Média, obtida nas região Centro-Oeste ( 41,2 ), e a Média da região Sudeste ( 51,3 ).
Não há registro de participante de curso a Distância de Bacharelado.
Isso significa que os valores tabelados para os cursos de Educação Presencial coincidem com os valores do Brasil, como um todo.
A Média das notas dos alunos de cursos de Bacharelado, foi 54,1, sendo que os alunos da região Norte obtiveram a Média mais baixa ( 50,2 ), e os da região Sul obtiveram a Média mais alta ( 55,9 ).
As demais médias foram : 54,7 na região Nordeste, 54,0 na região Sudeste e 51,7 na região Centro-Oeste.
A região que obteve a maior nota Máxima foi a Sul ( 98,6 ), e a com menor nota Máxima foi a Norte ( 89,8 ).
A Mediana do Brasil, como um todo, foi 55,5, sendo a maior Mediana obtida na região Sul ( 57,8 ), e a menor obtida na Norte ( 52,0 ).
A nota Mínima foi zero em todas as regiões.
Considerando -se as notas, segundo Grande Região, observa -se que existem diferenças estatisticamente significativas ao nível de 95 \% entre as duas menores médias, das regiões Norte e Centro-Oeste, e as médias das outras três regiões.
Entre as três maiores médias, não há diferença estatisticamente significativa.
O Gráfico 6.2 propicia a avaliação do desempenho dos estudantes do Componente de Formação Geral, a partir do histograma da distribuição das notas correspondentes.
A distribuição é unimodal, com moda em ( 50 ; 60 ] para o Bacharelado, intervalo superior ao da prova como um todo.
A moda para a Licenciatura foi o intervalo ( 40 ; 50 ], mesmo intervalo modal da distribuição de notas da prova, como um todo, para os alunos de Licenciatura ( Gráfico 6.1 ).
Nota -se, ainda, que no Gráfico 6.2, as notas apresentam dispersão maior, quando comparadas às do Gráfico 6.1 ( distribuição das notas da prova ), confirmado pela comparação dos desvios padrões ( 13,3 e 13,6 na prova como um todo, respectivamente para a Licenciatura e o Bacharelado, e, na mesma ordem, 17,5 e 17,9 na parte de Formação Geral ).
Na Tabela 6.4, são apresentadas as informações referentes ao desempenho dos concluintes do Componente de Formação Geral, em diferentes agregações : Categoria Administrativa, Organização Acadêmica e Habilitação.
Primeiramente, com respeito aos alunos de Licenciatura, observa -se que não existe diferença estatisticamente significativa entre as médias dos tipos de Categoria Administrativa : a maior Média obtida por estudantes de IES Públicas ( 49,5 ), e a menor, pelos de IES Privadas ( 50,3 ).
A ordenação decrescente das médias por Organização Acadêmica foi : Centros Universitários ( 53,6 ), Universidades ( 49,6 ), CEFET/IFET ( 48,9 ) e Faculdades ( 45,9 ).
Observa -se que não existe diferença estatisticamente significativa ao nível de 95 \% entre a maior Média da Organização Acadêmica Centro Universitário ( 53,6 ) e as demais Organizações Acadêmicas.
Entre os alunos de Bacharelado, a situação das médias com relação à Categoria Administrativa se inverte : maior para IES Públicas ( 57,2 ) do que para IES Privadas ( 51,5 ), com diferença estatisticamente significativa.
As Médias dos cursos de Bacharelado foram : CEFET/IFET ( 55,6 ), Universidades ( 55,2 ), Centros Universitários ( 52,9 ) e Faculdades ( 48,6 ), e há diferença estatisticamente significativa entre a menor média, obtida nas Faculdades ( 48,6 ) e as médias dos demais tipos de Organização Acadêmica. 6.1.3 Estatísticas Básicas do Componente de Conhecimento Específico A Tabela 6.5 apresenta as Estatísticas Básicas referentes ao Componente de Conhecimento Específico da área de Ciência da Computação ( Licenciatura/Bacharelado ) por Grande Região, segundo Habilitação e Modalidade de Ensino.
Quanto ao desempenho de Conhecimento Específico dos alunos dos cursos de Licenciatura, a Média foi 37,8.
A maior Média foi obtida na região Sul ( 45,0 ), e a menor, na região Centro-Oeste ( 32,6 ).
As demais médias foram : 34,3 na região Norte, 39,2 na região Sudeste e 39,5 na região Nordeste.
Já para os alunos dos cursos presenciais, a Média do desempenho de Conhecimento Específico foi maior, 38,2.
A maior Média foi obtida na região Sul ( 45,0 ), e a menor, na região Centro-Oeste ( 33,6 ).
As demais médias foram : 40,5 na região Nordeste, 40,4 na região Sudeste e 34,7 na região Nordeste.
Quanto à variabilidade das notas, o Desvio padrão do Brasil, como um todo, foi 14,4, sendo o maior Desvio padrão observado na região Centro-Oeste ( 16,8 ), e o menor, na região Norte ( 12,2 ).
Os demais desvios foram : 13,0 na região Sul, 13,4 na Nordeste e 15,0 na região Sudeste.
A Mediana das notas dos estudantes de cursos presenciais de Licenciatura de todo o Brasil foi 38,2.
A maior Mediana ocorreu na região Sul ( 44,4 ), e a menor, na região Centro-Oeste ( 29,8 ).
As demais medianas foram : 34,7 na região Norte, 40,6 na região Nordeste e 41,3 na região Sudeste.
A nota Máxima do Brasil, como um todo, foi 83,0, sendo obtida por, pelo menos, um aluno na região Sul.
As demais notas máximas foram : 78,7 na região Nordeste, 76,3 na região Centro-Oeste, 75,2 na região Norte e 71,8 na região Sudeste.
A nota Mínima foi zero, obtida no Centro-Oeste.
As demais notas mínimas foram : 3,7 na região Nordeste, 4,0 na Norte, 11,1 no Sudeste e 15,5 na região Sul.
Observa -se que existe diferença estatisticamente significativa da menor Média das notas do Componente de Conhecimento Específico, obtida na região Centro-Oeste ( 33,6 ), em relação às duas maiores, regiões Nordeste ( 40,5 ) e Sul ( 45,0 ).
Já a Média do desempenho de Conhecimento Específico dos alunos dos cursos a Distância de Licenciatura foi um pouco menor, 37,0.
A maior Média foi obtida na região Sul ( 45,3 ), e a menor, na região Norte ( 26,1 ).
As demais médias foram : 37,8 na região Nordeste, 38,6 na região Sudeste e 28,8 na região Centro-Oeste.
O Desvio padrão de todos os alunos de cursos a Distância de Licenciatura foi 14,1, sendo o maior Desvio padrão observado na região Sudeste ( 15,6 ), e o menor, na região Sul ( 8,1 ).
Os demais desvios foram : 9,2 na região Norte, 12,6 na região Nordeste e 8,6 na Centro-Oeste.
A Mediana das notas dos estudantes de cursos a Distância de Licenciatura foi 36,4.
A maior Mediana ocorreu na região Sul ( 45,3 ), e a menor, na região Norte ( 25,8 ).
No Nordeste a Mediana foi 37,7 no Sudeste, 37,0 e 29,6 na região Centro-Oeste.
A nota Máxima do Brasil, como um todo, foi 80,0, sendo obtida por, pelo menos, um aluno na região Sudeste.
As demais notas máximas foram : 38,5 na região Norte, 47,9 na região Centro-Oeste, 51,0 na região Sul e 65,8 na região Nordeste.
A nota Mínima foi zero, obtida na região Nordeste.
As notas mínimas das outras regiões foram : 11,1 na Norte, 3,7 na Sudeste, 39,5 na Sul e 11,1 na região Centro-Oeste.
Observa -se que existe diferença estatisticamente significativa entre a Média das notas do Componente de Conhecimento Específico da região Norte, a menor das médias, e a Média das notas das regiões Nordeste e Sudeste.
Como já foi mencionado, não houve participante em Educação a Distância de cursos de Bacharelado.
Assim sendo, os registros nas categorias Educação Presencial e Total coincidem.
Para os alunos de cursos de Bacharelado, a Média do desempenho de Conhecimento Específico foi 39,6.
A maior Média foi obtida na região Sul ( 43,0 ), e a menor, na região Norte ( 34,5 ).
As demais médias foram : 40,4 na região Nordeste, 38,8 na região Sudeste e 37,8 na região Centro-Oeste.
O Desvio padrão dos alunos de Bacharelado foi 15,1, sendo o maior desvio e A Mediana das notas dos estudantes de cursos de Bacharelado foi 38,4.
A maior Mediana ocorreu na região Sul ( 42,4 ), e a menor, na região Norte ( 32,4 ).
As demais medianas foram : 38,9 na região Nordeste, 36,9 na região Sudeste e 37,4 na região Centro-Oeste.
A nota Máxima do Brasil, como um todo, foi 99,3, obtida por pelo menos, um aluno na região Sudeste.
As demais notas máximas foram : 80,8 na região Norte, 92,9 na região Nordeste, 87,9 na região Sul e 79,1 na região Centro-Oeste.
A nota Mínima foi zero em todas as regiões, com exceção da região Norte com 3,6.
Observa -se que, para os alunos de cursos de Bacharelado, existem diferenças estatisticamente significativas entre as médias das notas do Componente de Conhecimento Específico de todas as Grandes Regiões, com exceção das regiões Sudeste e Centro-Oeste.
Assim como os Gráficos 6.1 e 6.2, o Gráfico 6.3, apresentado a seguir, permite uma avaliação do desempenho de concluintes de Ciência da Computação ( Licenciatura/Bacharelado ) segundo a Habilitação, em relação ao Componente de Conhecimento Específico com um histograma da distribuição das notas correspondentes.
Esta também é uma distribuição unimodal para ambas as habilitações, e o grupo modal é o mesmo para estudantes de Licenciatura e Bacharelado : ( 30 ; 40 ].
A Tabela 6.6 apresenta uma comparação dos resultados em relação à Categoria Administrativa e à Organização Acadêmica, segundo a Habilitação, agora levando -se em conta o desempenho de alunos do Componente de Conhecimento Específico da prova.
No que se refere à Organização Acadêmica, para os alunos de Licenciatura, a maior Média foi a dos Centros Universitários ( 43,0 ), vindo a seguir a de Universidades ( 38,5 ), a dos CEFET/IFET ( 36,0 ), e, por último, a das Faculdades ( 29,8 ).
O maior Desvio padrão, acima do valor para o Brasil, como um todo, ( 14,3 ), foi o das Faculdades ( 15,9 ).
As Universidades obtiveram a maior nota Máxima, igual a 83,0, com os Centros Universitários em segundo lugar, com 80,0.
As Faculdades obtiveram nota Máxima igual a 71,8, e os CEFET/IFET, 70,6.
As medianas foram : 42,9 nos Centros Universitários, a maior Mediana, 38,2 nas Universidades, 35,7 nos CEFET/IFET, e 25,8 nas Faculdades, a menor Mediana.
A nota Mínima foi zero para Universidades e Faculdades.
Quanto à Categoria Administrativa, observa -se um comportamento semelhante àquele da parte de Formação Geral e à prova como um todo, ou seja, a Média das IES Públicas ( 37,9 ) é superior à Média das IES Privadas ( 37,5 ), mas para esta componente tampouco, a diferença entre elas é estatisticamente significativa.
Para os alunos de Bacharelado, a maior Média, considerando -se a Organização Acadêmica, foi a de Universidades ( 41,8 ), e a menor, a das Faculdades ( 30,7 ).
As Universidades obtiveram a maior nota Máxima, igual a 99,3.
As medianas foram : 40,4 nas Universidades, 34,6 nos Centros Universitários, 29,3 nas Faculdades e 41,6 nos CEFET/IFET.
A nota Mínima foi zero em todas as Organizações Acadêmicas.
Quanto à Categoria Administrativa, observa -se um comportamento semelhante àquele da parte de Formação Geral e à prova, como um todo, para os alunos de Bacharelado, ou seja, nesse caso também a maior Média foi obtida por alunos de IES Públicas de ensino ( 45,4 ), com diferença estatisticamente significativa em relação à Média das IES Privadas ( 34,5 ). 6.2 ANÁLISE DAS QUESTÕES OBJETIVAS Esta seção apresenta estatísticas selecionadas e histogramas das Questões Objetivas de Formação Geral para Licenciatura ( 6.2.1 ) e para Bacharelado ( 6.2.2 ), e das Questões Objetivas de Conhecimento Específico para Licenciatura ( 6.2.3 ) e para Bacharelado ( 6.2.4 ).
São também apresentadas e comparadas as médias das subpopulações caracterizadas por Grande Região. 6.2.1 Componente de Formação Geral – Ciência da Computação ( Licenciatura ) A Tabela 6.7a apresenta as Estatísticas Básicas relativas às oito questões objetivas do componente da prova, que abrange a Formação Geral dos estudantes de Ciência da Computação ( Licenciatura ), segundo a Modalidade de ensino.
A Média dos alunos nestas questões foi 48,2.
A menor Média foi encontrada na região Norte ( 44,0 ), e a maior, na região Sul ( 52,6 ).
As demais médias foram : 46,4 na região Centro-Oeste, 49,6 na região Nordeste e 49,2 na região Sudeste.
Para os alunos de cursos Presenciais, a Média também foi 48,2.
A menor Média foi encontrada na região Norte ( 44,0 ), e a maior, na região Sul ( 52,2 ).
As demais médias foram : 45,7 na região Sudeste, 51,0 na região Nordeste e 48,1 na região Centro-Oeste.
Para os cursos Presenciais, o Desvio padrão do Brasil foi 20,7, sendo o maior Desvio padrão encontrado na região Centro-Oeste ( 23,1 ), e o menor, na região Sul ( 17,8 ).
Os demais desvios foram : 20,1 na região Norte, 19,8 na região Nordeste e 21,1 na região Sudeste.
Ainda em relação às Licenciatura Presenciais, as medianas do Brasil como um todo e em quatro das cinco regiões foram iguais a 50,0, na região Norte a Mediana foi 37,5.
A nota Máxima 100,0 foi alcançada nas regiões Norte, Nordeste e Centro-Oeste.
Nas regiões Sul e Sudeste a nota Máxima foi 87,5.
A nota mínima ( zero ) foi obtida por, pelo menos, um estudante em todas as regiões.
Para os alunos de cursos a Distância, a Média também foi 48,2.
A menor Média foi encontrada na região Centro-Oeste ( 39,9 ), e a maior, na região Sul ( 68,8 ).
As demais médias foram : 43,1 na região Norte, 47,2 na região Nordeste e.50,9 na região Sudeste.
O Desvio padrão do Brasil para esta Modalidade de ensino nos cursos de Licenciatura foi 22,3, sendo o maior Desvio padrão encontrado na região Sul ( 26,5 ), e o menor, na região Norte ( 15,5 ).
Os demais desvios foram : 23,0 na região Nordeste, 22,0 na região Sudeste e 19,9 na região Centro-Oeste.
As medianas do Brasil, como um todo, e para as regiões Norte, Nordeste e Sudeste foi 50,0.
Foi mais baixa para as região Centro-Oeste ( 37,5 ) e mais alta para a região Sul ( 68,8 ).
A nota Máxima, 100,0, foi alcançada nas regiões Nordeste e Sudeste, e a Mínima ( zero ), nas regiões Nordeste, Sudeste e Centro-Oeste.
A Tabela 6.8a apresenta o Índice de Facilidade e o Índice de Discriminação ( ponto-bisserial ) para cada uma das questões objetivas do Componente de Formação Geral.
Quanto ao Índice de Facilidade, foram usadas as seguintes cores para diferenciar o nível de dificuldade da questão :  Azul para as questões classificadas com índice Muito fácil ( \&gt ; =0,86 ), verde para as questões classificadas com índice Fácil ( 0,61 a 0,85 ), amarelo para as questões classificadas com Médio ( 0,41 a 0,60 ), vermelho para as questões classificadas como Difícil ( 0,16 a 0,40 ) e roxo para as questões classificadas como Muito difícil ( \&lt ; =0,15 ).
Já quanto ao Índice de Discriminação, foram usadas as seguintes cores para qualificar a questão :  As questões classificadas com índice Fraco receberam a cor vermelha ( \&lt ; =0,19 ), as classificadas com Médio receberam a cor amarela ( 0,20 a 0,29 ), as classificadas com Bom receberam a cor verde ( 0,30 a 0,39 ), e as classificadas com Muito bom ( \&gt ; =0,40 ), a cor azul.
As questões objetivas do Componente de Formação Geral, segundo o Índice de Facilidade, foram assim avaliadas para os estudantes de Licenciatura : das oito questões, nenhuma teve o Índice de Facilidade classificado como Muito fácil e uma questão foi tida como Fácil.
Quatro questões foram consideradas com índice Médio, situando -se no intervalo entre 0,41 e 0,60 do Índice de Facilidade, ou seja, houve entre 41,0 \% e 60,0 \% de acertos, enquanto três questões foram classificadas na categoria Difícil, situando -se no intervalo entre 0,16 e 0,40.
Nenhuma questão foi classificada como Muito difícil.
O Índice de Facilidade variou de 0,20 a 0,80, e o de Discriminação, de 0,36 a 0,52.
A Tabela 6.9a apresenta a distribuição das questões, considerando -se simultaneamente a classificação dos dois índices.
As sete questões com Índice de Discriminação Muito bom figuraram entre três níveis de dificuldade Fácil, Médio e Difícil : uma classificada na categoria Fácil ( questão 6 ) do Índice de Facilidade, quatro na categoria Médio ( questões 2, 4, 5 e 8 ) e duas na categoria Difícil ( questões 1 e 7 ).
Em particular, a questão 4 apresentou o maior poder discriminatório, com índice 0,52, e seu Índice de Facilidade foi Médio com 59 \% dos estudantes acertando -a.
A questão 6 obteve o máximo de acertos, com uma proporção de 0,80 acertos e seu Índice de Discriminação foi Muito Bom ( 0,51 ).
A questão de número 3 apresentou um Índice de Facilidade de 0,35 ( Difícil ), ou seja, 35 \% dos estudantes conseguiu resolvê -la, dentro do universo de participantes, e teve Índice de Discriminação Bom ( 0,36 ).
O Gráfico 6.4a, para exemplificar, analisa o comportamento da questão de número 4 de Formação Geral.
Trata -se de uma questão considerada Médio, com relação à facilidade, e uma das duas que obteve o maior Índice de Discriminação dessa parte da prova ( 0,52 ).
Neste gráfico, cada uma das cinco curvas representa o percentual de respostas em determinada alternativa da questão, em função do número de acertos dos estudantes nessa parte da prova ( Formação Geral/Múltipla Escolha ), antes de possíveis eliminações pelo critério do ponto-bisserial.
A curva em verde corresponde à alternativa B, a alternativa correta para essa questão.
Observa -se que entre os estudantes com menor número de acertos, nessa parte do exame, a situação mais frequente foi a escolha da alternativa E ( em vermelho ), incorreta.
Por exemplo, entre os estudantes que acertaram três questões, 48,2 \% escolheu a alternativa correta B ( em verde ), cerca de 23,4 \% escolheram a alternativa E ( em vermelho ), 17,1 \% escolheram a alternativa D ( em roxo ), 7,7 \% escolheram a alternativa A ( em azul ), e 3,2 \%, a C ( em preto ).
Entre os que acertaram três respostas entre as questões de múltipla escolha de Formação Geral, 0,0 \% deixou a questão em branco e 0,5 \% marcou mais de uma alternativa, invalidando a questão.
À medida em que o número de acertos aumenta, indicando desempenho melhor nessa parte da prova, aumenta concomitantemente a proporção de estudantes que selecionaram a alternativa correta B, atingindo 100 \% para os estudantes com oito acertos.
Essa análise permite verificar como a questão discriminou os grupos de desempenho, justificando o alto índice obtido na questão.
Cumpre notar que não é possível inferir deste gráfico nem o índice de facilidade ( que seria uma média da proporção ponderada pela quantidade de alunos com cada uma das notas ), nem o índice de discriminação ponto-bisserial, por razão equivalente.
No caso extremo, no qual a grande concentração dos acertos dos alunos fosse abaixo de 4, o índice de facilidade seria obrigatoriamente abaixo de 40 \% ( nesse exemplo ).
Caso a concentração fosse em seis acertos ou mais, o índice seria obrigatoriamente acima de 60 \%.
Os gráficos relativos às demais questões de Formação Geral constam do Anexo Ia. 6.2.2 Componente de Formação Geral – Ciência da Computação ( Bacharelado ) A Tabela 6.7b apresenta as Estatísticas Básicas relativas às oito questões objetivas do componente, que abrange a Formação Geral dos estudantes de Ciência da Computação ( Bacharelado ), segundo a Modalidade de ensino.
Cabe destacar que não houve participante de cursos de Educação a Distância.
Portanto, os valores registrados para Educação Presencial coincidem com os valores classificados como Total.
A Média dos alunos de Bacharelado nestas questões foi 55,5.
A maior Média foi encontrada na região Sul ( 58,8 ), e a menor, na região Norte ( 50,0 ).
As demais médias foram : 56,1 na região Nordeste, 54,7 na região Sudeste e 53,9 na região Centro-Oeste.
O Desvio padrão do Brasil foi 21,4, sendo o maior Desvio padrão encontrado na região Centro-Oeste ( 22,2 ), e o menor, na região Sul ( 20,9 ).
Os desvios das demais Grandes Regiões foram bem próximos : 21,5 nas regiões Norte e Nordeste e 21,2 na região Sudeste.
A Mediana do Brasil como um todo e de duas das cinco regiões ( Sudeste e Sul ) foi 62,5.
A Mediana nas regiões Norte, Sudeste e Sul foi 50,0.
A nota Máxima 100,0 foi alcançada em todas as Grandes Regiões.
As notas mínimas ( zero ) foram iguais para todas as regiões.
A Tabela 6.8b apresenta o Índice de Facilidade e o Índice de Discriminação ( ponto-bisserial ) para cada uma das questões objetivas do Componente de Formação Geral.
A convenção das cores é a mesma da Tabela 6.8a.
As questões objetivas do Componente de Formação Geral, segundo o Índice de Facilidade, foram assim avaliadas : das oito questões, uma teve o Índice de Facilidade classificado como Muito fácil, e duas questões foram tidas como Fácil.
Três questões foram consideradas com índice de dificuldade Médio, enquanto duas questões foram classificadas na categoria Difícil, situando -se no intervalo entre 0,16 e 0,40.
Nenhuma das questões apresentou menos de 15 \% de acertos, e, portanto, não houve questão classificada como Muito difícil.
O Índice de Facilidade variou de 0,34 a 0,88, e o de Discriminação, de 0,39 a 0,51.
A Tabela 6.9b apresenta a distribuição das questões, considerando -se simultaneamente a classificação dos dois índices, semelhantemente ao apresentado na Tabela 6.9a, mas aqui para os alunos de Bacharelado.
As sete questões com Índice de Discriminação Muito bom figuraram entre os quatro níveis de dificuldade Muito fácil, Fácil Médio e Difícil : uma classificada na categoria Muito fácil ( a questão 6 ) do Índice de Facilidade, duas na categoria Fácil ( questões 4 e 5 ), duas na categoria Médio ( questões 3 e 8 ) e, por fim, duas na categoria Difícil ( questões 1 e 7 ).
As questões 4 e 5 apresentaram o maior poder discriminatório, com índice 0,51, e foram consideradas Fácil, em termos de facilidade, com uma proporção de 0,73 e 0,67 acertos, respectivamente.
O máximo de acertos, dentre as oito questões objetivas de Formação Geral, foi alcançado na questão 6, com 88 \% dos participantes marcando a opção correta.
O Gráfico 6.4b, para exemplificar, analisa o comportamento da questão de número 4 de Formação Geral.
Trata -se de uma questão considerada Fácil, com relação à facilidade, e foi a com Índice de Discriminação Muito bom.
Semelhantemente ao apresentado no Gráfico 6.4a, neste gráfico também, cada uma das cinco curvas representa o percentual de respostas em determinada alternativa da questão, em função do número de acertos dos estudantes nessa parte da prova ( Formação Geral/Múltipla Escolha ), antes de possíveis eliminações pelo critério do ponto-bisserial.
A curva em verde corresponde à alternativa B, a alternativa correta para essa questão.
Observa -se que entre os estudantes com menor número de acertos, nessa parte do exame, a situação mais frequente foi a escolha das alternativas D ( em roxo ) e E ( em vermelho ), incorretas.
Por exemplo, entre os estudantes que acertaram três questões, mais da metade ( 56,0 \% ) escolheram a alternativa correta B ( em verde),17,3 \% escolheram a alternativa D ( roxo ), 15,2 \% escolheram a E ( em vermelho ), 8,1 \% escolheram a alternativa A ( azul ), e 3,2 \%, as alternativas C ( em preto ).
Entre os que acertaram três respostas entre as questões de múltipla escolha de Formação Geral, 0,1 \% deixou a questão em branco e 0,1 \% marcou mais de uma alternativa, invalidando a questão.
À medida que o número de acertos aumenta, indicando desempenho melhor nessa parte da prova, aumenta concomitantemente a proporção de estudantes que selecionaram a alternativa correta B, atingindo 100 \% para os estudantes, com oito acertos.
Essa análise permite verificar como a questão discriminou os grupos de desempenho, justificando o alto índice obtido na questão.
Os gráficos relativos às demais questões de Formação Geral constam do Anexo Ib. 6.2.3 Componente de Conhecimento Específico – Ciência da Computação ( Licenciatura ) A Tabela 6.10a apresenta as Estatísticas Básicas em relação às questões objetivas do Componente de Conhecimento Específico da prova de Ciência da Computação ( Licenciatura ), por Grande Região, segundo a Modalidade de ensino.
Nessa parte da prova, nenhuma das questões foi anulada pela Comissão Assessora da Área.
Como será discutido ainda nesta seção, quatro questões objetivas do Componente de Conhecimentos Específicos não foram usadas no cômputo das notas por terem sido descartadas pelo critério do ponto-bisserial.
Assim, as notas foram calculadas com base em 23 das 27 questões objetivas de Conhecimentos Específicos.
A Média deste componente para os alunos de cursos de Licenciatura foi de 41,6.
A menor Média foi observada na região Centro-Oeste ( 36,2 ), e a maior, na região Sul ( 48,4 ).
O Desvio padrão de todo o Brasil foi 15,6, sendo o menor Desvio padrão encontrado na região Norte ( 13,5 ), e o maior, na região Sudeste ( 16,5 ).
A situação é um pouco diferente quando se desagrega por Modalidade de Ensino : a Média desse componente para os alunos de cursos Presenciais de Licenciatura foi de 41,9.
A menor Média foi observada na região Norte ( 38,3 ), e a maior, na região Sul ( 48,4 ).
O Desvio padrão de todo o Brasil foi 15,7, sendo o menor Desvio padrão encontrado na região Norte ( 13,5 ), e o maior, na região Centro-Oeste ( 18,1 ).
A Mediana de todo o Brasil foi 43,5, a mesma das regiões Nordeste e Sudeste.
Nas regiões Norte ( 39,1 ) e Centro-Oeste ( 34,8 ), a Mediana foi menor que a nacional, e na região Sul ( 47,8 ), foi maior.
A nota Máxima da prova foi 87,0, obtida nas questões objetivas do Componente de Conhecimento Específico, por, pelo menos, um aluno de Licenciatura Presencial das regiões Nordeste e Sul.
Nas demais regiões as notas Máximas foram : 82,6 na região Norte, 73,9 no Sudeste e 82,6 no Centro-Oeste.
A nota Mínima foi zero nas regiões Norte e Centro-Oeste.
Na região Nordeste, a nota Mínima foi 4,3, na região Sudeste foi 13,0 e na região Sul, 17,4.
Já para os alunos de cursos a Distância de Licenciatura, a Média desse componente foi mais baixa, 40,9.
A menor Média foi observada na região Norte ( 28,5 ), e a maior, na região Sul ( 50,0 ).
O Desvio padrão de todo o Brasil foi 15,3, sendo o menor Desvio padrão encontrado na região Sul ( 9,2 ), e o maior, na região Sudeste ( 16,8 ).
A Mediana de todo o Brasil foi 39,1, a mesma da região Sudeste.
Na região Norte a Mediana foi 30,4, na Nordeste foi 43,5, na Centro-Oeste 34,8 e na região Sul, a Mediana foi maior ( 50,0 ).
A nota Máxima da prova foi 82,6, obtida nas questões objetivas do Componente de Conhecimento Específico, por, pelo menos, um aluno da região Sudeste.
A nota Mínima ( zero ) foi obtida na região Nordeste.
A Tabela 6.11a apresenta os Índices de Facilidade e Discriminação ( ponto-bisserial ) das questões objetivas do Componente de Conhecimento Específico para os estudantes de Ciência da Computação ( Licenciatura ).
Para facilitar a diferenciação das questões, serão usadas as mesmas cores da Tabela 6.8a e 6.8b para as diferentes classificações dos Índices de Facilidade e de Discriminação.
Observa -se que nenhuma das questões foi anulada pela CAA, e que as análises dizem respeito às 27 questões válidas.
A partir do Índice de Facilidade obtido, pode -se concluir que mais do que metade das questões objetivas da prova foi considerada, pelo menos, Difícil : das 27 questões, 13 foram classificadas como Difícil, e uma, como Muito difícil.
Três questões foram classificadas como Fácil, e outras dez foram consideradas como Médio.
Já quanto ao Índices de Discriminação das questões objetivas do Componente de Conhecimento Específico da prova, obteve -se como resultado a seguinte classificação : 11 das 27 questões foram consideradas como boas, enquanto quatro delas tiveram Índice de Discriminação Muito bom.
Assim, para um pouco mais do que a metade das questões, – 15 em 27, – os Índices de Discriminação foram Bom ou Muito bom.
Dentre as demais, oito delas foram classificadas como Médio, e outras quatro, como Fraco, sendo 12, por conseguinte, a quantidade de questões nos dois patamares mais baixos de discriminação.
Constata -se, assim, que a prova, – no que se refere ao Componente de Conhecimento Específico, – possuía média capacidade de discriminar dentre aqueles que dominam ou não o conteúdo.
A Tabela 6.12a apresenta a distribuição das questões, considerando -se simultaneamente a classificação dos dois índices.
Dentre as questões que alcançaram os maiores índices de discriminação, classificadas com Índice de Discriminação Muito bom, uma delas, a de número 22, obteve o índice mais elevado ( 0,49 ) e foi classificada como Médio quanto ao Índice de Facilidade, com 60 \% dos estudantes, marcando a opção correta.
A questão de número 10 foi a mais difícil dentre as 27 questões específicas, com baixo Índice de Facilidade, apenas 10 \% de acertos.
Essa questão apresentou poder discriminatório baixo, 0,09, o que comprova ter sido a mais difícil para os estudantes.
Destaca -se, também, a questão 21, com Índice de Facilidade 0,17, o que, em termos percentuais, corresponde a 17 \% de estudantes, que responderam acertadamente e com Índice de Discriminação de 0,08.
Tais questões foram, portanto, pelo critério ponto-bisserial, consideradas inadequadas.
Por isso, as questões 9, 10 e 21, além da questão 27, foram eliminadas do cômputo da nota final.
A título de exemplo das análises do comportamento das questões objetivas, o Gráfico 6.5a analisa a questão 22 do Componente de Conhecimento Específico.
Esta questão foi considerada pelas respostas dos estudantes avaliados na prova como uma questão com nível de facilidade Médio, 0,60, ou seja, 60 \% dos estudantes assinalaram acertadamente a opção E, correspondente ao gabarito.
Como já comentado, seu Índice de Discriminação foi igual a 0,49, classificado como Muito bom, o maior valor de discriminação.
Neste gráfico também, cada uma das cinco curvas representa o percentual de respostas em determinada alternativa da questão 22, em função do número de acertos dos estudantes nessa parte da prova, antes de possíveis eliminações de questões pelo critério do ponto-bisserial.
A alternativa correta E, representada no gráfico pela curva em vermelha, foi escolhida em maiores proporções pelos alunos com desempenho melhor nessa parte da prova.
Já as alternativas incorretas, também denominadas distratores, foram selecionadas, principalmente, por aqueles com notas mais baixas.
Observa -se que a soma não é 100 \%, por causa das questões não respondidas ou com mais de uma opção marcada.
Aqueles com nota zero, na sua quase totalidade, deixaram essa questão em branco ou marcaram mais de uma alternativa, comportamento considerado inválido.
A proporção de alunos que selecionou a resposta correta E aumenta gradativamente, chegando a atingir 100 \% para 17 acertos ou mais, com uma descontinuidade em 21 acertos, enquanto a proporção dos que escolheram alternativas incorretas decai, a partir de seis ou sete acertos, como função do número de acertos nessa parte da prova.
Os gráficos relativos às demais questões do Conhecimento Específico constam do Anexo I. 6.2.4 Componente de Conhecimento Específico – Ciência da Computação ( Bacharelado ) A Tabela 6.10b apresenta as Estatísticas Básicas em relação às questões objetivas do Componente de Conhecimento Específico da prova de Ciência da Computação ( Bacharelado ) por Grande Região segundo a Modalidade de Ensino.
Nenhuma das questões foi anulada pela Comissão Assessora da Área.
Porém, como será discutido ainda nesta seção, três das questões objetivas do Componente de Conhecimentos Específicos não foram usadas no cômputo das notas, por terem sido descartadas pelo critério do ponto-bisserial.
Assim, as notas foram calculadas com base em 24 das 27 questões objetivas de Conhecimentos Específicos.
A Média deste componente para os alunos de cursos de Bacharelado foi de 41,6.
A menor Média foi observada na região Norte ( 37,0 ), e a maior, na região Sul ( 44,6 ).
O Desvio padrão de todo o Brasil foi 14,6, sendo o menor Desvio padrão encontrado na região Norte ( 13,2 ), e o maior, na região Centro-Oeste ( 14,9 ).
A Modalidade de Ensino Educação a Distância não apresentou participante nos cursos de Bacharelado.
Portanto, todos os valores referem -se a integrantes de Educação Presencial.
A Mediana de todo o Brasil foi 41,7, a mesma nas regiões Nordeste, Sudeste e Centro-Oeste.
Nas regiões Norte e Sul, as medianas foram 37,5 e 45,8, respectivamente.
A nota Máxima da prova ( 100 ) foi obtida nas questões objetivas do Componente de Conhecimento Específico por, pelo menos, um aluno da região Sudeste.
A nota Mínima foi zero em todas as regiões, com exceção da região Norte, cuja nota Mínima foi 4,2.
A Tabela 6.11b apresenta os Índices de Facilidade e Discriminação ( ponto-bisserial ) das questões objetivas do Componente de Conhecimento Específico para os estudantes de Ciência da Computação ( Bacharelado ).
Para facilitar a diferenciação das questões, serão usadas as mesmas convenções definidas para as Tabela 6.8a, 6.8b e 6.11a.
A partir do Índice de Facilidade obtido, pode -se concluir que dois terços das questões objetivas da prova fora considerados Difícil : das 27 questões válidas, 18 foram classificadas como Difícil e nenhuma como Muito difícil.
Quatro questões foram classificadas como Fácil, e outras cinco, consideradas como Médio.
Já quanto ao Índice de Discriminação das questões objetivas do Componente de Conhecimento Específico da prova, obteve -se como resultado a seguinte classificação : 12 das 27 questões foram consideradas como boas, enquanto outras duas tiveram Índice de Discriminação Muito bom.
Assim, para um pouco mais da metade das questões – 14 em 27 – os Índices de Discriminação foram Bom ou Muito bom.
Dentre as demais, dez delas foram classificadas como Médio e três como Fraco, sendo 13, por conseguinte, a quantidade de questões nos dois patamares mais baixos de discriminação.
Constata -se, assim, que a prova, – no que se refere ao Componente de Conhecimento Específico – possuía média capacidade de discriminar dentre aqueles que dominam ou não o conteúdo.
A Tabela 6.12b apresenta a distribuição das questões, levando -se em conta simultaneamente a classificação dos dois índices.
Dentre as questões que alcançaram os maiores índices de discriminação, classificadas com Índice de Discriminação Muito bom, uma delas, a de número 11, teve o índice mais elevado ( 0,42 ), sendo classificada como Médio quanto ao Índice de Facilidade, com 50 \% dos estudantes marcando a opção correta.
A questão de número 21 foi a mais difícil dentre as 27 questões específicas, com baixo Índice de Facilidade, apenas 18 \% de acertos.
Essa questão apresentou poder discriminatório muito baixo, 0,11, o que comprova ter sido a mais difícil para os estudantes.
Tal questão foi, portanto, pelo critério ponto-bisserial, considerada inadequada.
Por isso, a questão 21 foi eliminada do cômputo da nota final, assim coma as questões 25 e 26.
A título de exemplo das análises do comportamento das questões objetivas, o Gráfico 6.5b analisa a questão 11 do Componente de Conhecimento Específico.
Esta questão foi considerada, pelas respostas dos estudantes avaliados na prova, como uma questão com Índice de Facilidade Médio, 0,50, ou seja, 50 \% dos estudantes assinalaram acertadamente a opção E, correspondente ao gabarito.
Como já comentado, seu Índice de Discriminação foi igual a 0,42, classificada como Muito bom, o maior valor de discriminação.
Neste gráfico, cada uma das cinco curvas representa o percentual de respostas em determinada alternativa da questão 11, em função do número de acertos dos estudantes nessa parte da prova, antes de possíveis eliminações de questões pelo critério do ponto-bisserial.
A alternativa correta E, representada no gráfico pela curva em vermelho, foi escolhida em maiores proporções pelos alunos com desempenho melhor nessa parte da prova.
Já as alternativas incorretas, também denominadas distratores, foram selecionadas, principalmente, por aqueles com notas mais baixas.
Observa -se que a soma não é 100 \%, por causa das questões não respondidas ou com mais de uma opção marcada.
Aqueles com nota zero, na sua quase totalidade, deixaram essa questão em branco ou marcaram mais de uma alternativa, comportamento considerado inválido.
A proporção de alunos que selecionou a resposta correta E aumenta gradativamente, chegando a atingir 100 \% para 21 acertos, com uma queda entre aqueles que tiveram 23 acertos e um retorno a 100 \% para 24 acertos ou mais.
A proporção dos que escolheram alternativas incorretas decai, a partir de três ou quatro acertos, como função do número de acertos nessa parte da prova.
A exceção é a alternativa B ( em verde ) que apresenta um aumento com a proporção de acertos entre aqueles com pouco mais do que 20 acertos, um ponto atípico.
Os gráficos relativos às demais questões do Conhecimento Específico constam do Anexo I. 6.3 ANÁLISE DAS QUESTÕES DISCURSIVAS Esta seção apresenta estatísticas selecionadas e histogramas das Questões Discursivas de Formação Geral ( 6.3.1 ) e Conhecimento Específico ( 6.3.2 ).
São também apresentadas e comparadas as médias de subpopulações, caracterizadas por Grande Região. 6.3.1 Componente de Formação Geral As análises dos resultados de desempenho dos estudantes de Ciência da Computação ( Licenciatura/Bacharelado ) nas duas questões discursivas relativas à Formação Geral encontram -se na Tabela 6.13 e no Gráfico 6.6, desagregadas por Grande Região, segundo Habilitação e Modalidade de Ensino.
Na Tabela 6.13, observa -se que a nota Média nesse conjunto de questões foi acima da obtida nas objetivas pelos estudantes de Licenciatura.
Esses estudantes obtiveram, em Formação Geral, Média 48,2 nas questões objetivas, e 51,8 nas questões discursivas.
Já a situação \textbf{inversa} acontece para os alunos de Bacharelado : média maior nas objetivas ( 55,5 ) do que nas discursivas ( 52,2 ).
Considerando -se agora os alunos de cursos presenciais de Licenciatura, observa -se que a nota Média nesse conjunto de questões foi acima da obtida nas objetivas.
Esses estudantes obtiveram, em Formação Geral, Média 48,2 nas questões objetivas, e 51,9, nas questões discursivas.
Pode -se notar também que o Desvio padrão nesse conjunto de questões foi maior do que o obtido nas objetivas, 20,7 nas questões objetivas e 25,1 nas questões discursivas.
A maior Média foi obtida na região Sul ( 61,4 ), e a menor, na região Centro-Oeste ( 45,8 ).
A Mediana dos alunos de cursos presenciais de Licenciatura, neste componente, foi 58,0.
Nas regiões Norte e Nordeste ( 58,8 ), na região Sudeste ( 57,5 ) e Centro-Oeste ( 49,0 ) a menor Mediana e, na região Sul ( 66,5 ), foi maior que o valor para o Brasil.
A nota Máxima ( 96,0 ) foi a mesma nas regiões Norte e Sul, sendo 93,0 na região Nordeste, 90,0 na região Sudeste e 90,0 na região Centro-Oeste.
A nota Mínima ( zero ) foi a mesma em todas as regiões do Brasil.
Também para os alunos de cursos a Distância de Licenciatura a nota Média nesse conjunto de questões foi acima da obtida nas questões objetivas pelo mesmo grupo de estudantes.
Estes estudantes obtiveram, em Formação Geral, Média 48,2 nas questões objetivas e 51,5 nas questões discursivas.
Pode -se notar que o Desvio padrão nesse conjunto de questões também foi maior do que o obtido nas objetivas, 22,3 nas questões objetivas e 25,7 nas questões discursivas.
A maior Média foi obtida na região Sul ( 72,5 ), e a menor, na região Centro-Oeste ( 43,1 ).
A Mediana dos alunos de cursos a Distância de Licenciatura, neste componente, foi 57,0.
Na região Centro-Oeste ( 45,5 ) a Mediana foi menor e nas demais regiões foi maior que o valor para o Brasil : Norte ( 63,0 ), Nordeste e Sudeste ( 58,0 ) e Sul ( 72,5 ).
A nota Máxima ( 95,0 ) foi a mesma na região Nordeste, sendo 91,5 na região Norte, 94,5 na região Sudeste, 83,0 na Sul e 85,5 na região Centro-Oeste.
A nota Mínima ( zero ) foi a mesma nas regiões Nordeste, Sudeste e Centro-Oeste.
No Norte, a nota Mínima foi 35,0 e no Sul 62,0.
Já para os alunos do Bacharelado, a nota Média nas questões objetivas foi maior do que as das discursivas.
Esses estudantes obtiveram nas questões objetivas Média 55,5 e, nas questões discursivas, 52,2.
Pode -se notar que o Desvio padrão nesse conjunto de questões foi maior do que o obtido nas objetivas, 21,4 nas questões objetivas, e 25,9 nas questões discursivas.
A maior Média foi obtida na região Sudeste ( 52,9 ), e a menor, na região Centro-Oeste ( 48,3 ).
A Mediana dos alunos de cursos de Bacharelado, neste componente, foi 59,0, a mesma da região Sudeste.
Na região Norte ( 57,0 ), a região Nordeste com a Mediana mais alta ( 60,0 ), a região Sul ( 58,5 ) e a menor Mediana na região Centro-Oeste ( 56,0 ).
A nota Mínima ( zero ) foi a mesma em todas as regiões do Brasil.
Como já comentado, não há registro de participantes de curso a Distância de Bacharelado : os registros totais e os de Educação Presencial coincidem.
O Gráfico 6.6 representa a distribuição das notas nas questões discursivas do Componente de Formação Geral.
A moda desta distribuição ocorre no intervalo ( 60 ; 70 ] para os cursos de Licenciatura e de Bacharelado.
Destaca -se também o intervalo [ 0 ; 10 ] com uma proporção em torno de 12 \% do total de notas, para ambas as Habilitações, sendo que no intervalo [ 0 ; 10 ] inclui -se, além da nota zero, a frequência de alunos que deixaram a resposta a esse tipo de questão em branco.
Na sequência, os resultados verificados para cada uma das questões discursivas de Formação Geral serão apresentados, estabelecendo -se relações com os conteúdos abordados em cada uma delas.
Os comentários da Banca de docentes corretores, a respeito do observado na correção das respostas dos estudantes, suas impressões e conclusões serão apresentados junto à análise de cada questão.
Cumpre esclarecer que, tendo em vista que as questões discursivas de Formação Geral são padronizadas, ou seja, constam de todas as provas, os comentários da Banca são os mesmos para todas as carreiras acadêmicas, sendo direcionados a todos os estudantes que participaram do Enade/2017.
A seguir, serão analisados os desempenhos da Área de Ciência da Computação ( Licenciatura/Bacharelado ) nas duas questões discursivas de Formação Geral do Enade/2017, comparando -se os resultados obtidos com comentários para cada questão. 6.3.1.1 Análise de Conteúdo da Questão Discursiva 1 do Componente de Formação Geral Os dados de Ciência da Computação ( Licenciatura/Bacharelado ), obtidos a partir das respostas à questão 1, encontram -se na Tabela 6.14 e no Gráfico 6.7, desagregados por Habilitação e Modalidade de Ensino.
Nesta questão – de melhor desempenho dentre as duas de Formação Geral – os alunos, dos cursos de Licenciatura, tiveram Média 52,4.
A maior Média para a questão 1 entre os alunos de cursos de Licenciatura foi obtida na região Sul ( 61,2 ), e a menor, na região Centro-Oeste ( 45,0 ).
Os alunos dos cursos presenciais e a Distância de Licenciatura obtiveram médias mais altas na questão 1 : 52,4 para os alunos de cursos presenciais e 52,8 para os de cursos a Distância.
A maior Média para a questão 1 entre os alunos de cursos presenciais de Licenciatura foi obtida na região Sul ( 61,2 ), e a menor, na região Centro-Oeste ( 45,0 ).
Entre os alunos de cursos a Distância de Licenciatura, a maior Média foi obtida nas regiões Norte e Sul ( 70,0 ), e a menor, na região Centro-Oeste ( 47,3 ).
Quanto à variabilidade das notas, o Desvio padrão das notas foi 28,3, para os alunos de Educação Presencial, e 28,1, para os de Educação a Distância.
Para os cursos presenciais de Licenciatura, as medianas das regiões Norte e Sudeste e a do Brasil, como um todo, foram iguais ( 55,0 ).
Na região Nordeste, foi 60,0, na região Sul, 65,0 e na região Centro-Oeste, 50,0.
Já para os cursos a Distância, as medianas de todas as regiões e a do Brasil, com exceção da região Nordeste, foram mais altas.
As notas máximas da questão discursiva 1 foram, para ambas as Modalidades de Ensino, as mesmas para quase todas as regiões do Brasil : 100,0.
A Mínima foi zero em todas as regiões da Educação Presencial da Licenciatura.
A nota Mínima também foi zero nas regiões Nordeste, Sudeste e Centro-Oeste da Educação a Distância.
Entre os alunos de Bacharelado, a Média nesta questão foi ligeiramente maior do que na questão discursiva 2.
Os alunos dos cursos de Bacharelado obtiveram Média 52,1.
A maior Média para a questão 1, entre os alunos de cursos de Bacharelado, foi obtida na região Norte ( 53,7 ), e a menor, na região Centro-Oeste ( 49,0 ).
Quanto à variabilidade das notas, o Desvio padrão das notas foi 29,0, maior do que o Desvio Padrão da região Sudeste ( 28,3 ).
Os demais desvios padrões foram : 29,4 na região Norte, 29,6 na região Nordeste, 29,5 na região Sul, e 30,0, na Centro-Oeste.
Para os cursos de Bacharelado, a mediana do Brasil e das regiões Sudeste e Centro-Oeste foi 55,0.
Nas demais regiões a Mediana foi maior, 60,0.
As notas máximas e mínimas da questão discursiva 1 foram, sem exceção, as mesmas para todas as regiões do Brasil, respectivamente, 100,0 e 0,0.
Como já foi dito, não houve registro de participantes na Modalidade de Educação a Distância para o Bacharelado.
O Gráfico 6.7 mostra a distribuição das notas na questão discursiva 1 do Componente de Formação Geral.
Observa -se que a maior frequência corresponde ao intervalo ( 70 ; 80 ] para ambas as Habilitações.
Na Licenciatura, os intervalos ( 50, 60 ] e ( 60, 70 ] apresentam frequências bem próximas da frequência do intervalo modal.
Também é importante a proporção dos alunos que deixaram a resposta a essa questão em branco, entre 10 \% e 15 \% dos participantes em ambas as habilitações. 6.3.1.2 Comentários sobre a correção de Conteúdo das respostas à Questão Discursiva 1 O enunciado era claro e indicava um tema de extrema relevância, qual seja, a incidência de doença grave integrante do rol das DST ( Doenças Sexualmente Transmissíveis ), que encaminha para a análise do direito fundamental à saúde, que tem assento constitucional ( CRFB, art. 6º. caput : São direitos sociais a educação, a saúde, a alimentação, o trabalho, a moradia, o transporte, o lazer, a segurança, a previdência social, a proteção à maternidade e à infância, a assistência aos desamparados, na forma desta Constituição ).
A questão focalizava, por meio de três textos, um tema social extremamente relevante, qual seja, a relação entre homem e mulher nas suas interações íntimas, no caso representadas pela vida sexual, e o reflexo das relações de poder veiculadas pela sociedade, que apontam, no mais das vezes, o predomínio do homem.
Tal posição de superioridade é justificada pelas convenções arraigadas que acarretam a educação dos homens e mulheres por padrões patriarcais, que redundam em posições machistas.
Um dos textos demandava uma reflexão sobre a realidade dos casais que convivem em relações estáveis ou fugazes e não têm o controle dos efeitos das relações sexuais, que abrangem a gravidez não planejada e, em decorrência do não cuidado, a inserção de doenças transmitidas sexualmente.
Esse panorama é refletido na pouca influência das mulheres sobre os homens quanto à proteção, normalmente representada pelo uso de preservativos nas relações intimas, e após a constatação da doença, no encaminhamento das pessoas para tratamento.
Como efeito colateral, está a transmissão de doenças para os fetos.
Outro texto trazia à tona o maior cuidado que as mulheres têm, em geral, com sua saúde, notadamente no momento em que se descobre grávida e que redunda na realização de exames obrigatórios de pré-natal, que podem facilmente detectar doenças e, caso encontradas, ter o seu tratamento encaminhado.
A par disso, destaca -se que a crise de recursos financeiros agravou a situação, que já recebe influxos sociais negativos, como a dificuldade do homem em utilizar preservativos, realizar exames e procurar tratamentos quando adoece.
Como muito divulgado pelas mídias, a crise financeira do país repercutiu diretamente na prestação de serviços de saúde em todas as esferas : distrital, federal, estadual e municipal.
Assim, agravou -se o problema da falta de medicamentos essenciais para o tratamento da sífilis e, certamente, de outras doenças.
Como se verifica, a temática indicada no enunciado, representada pelos textos, pertence à realidade conhecida amplamente pelos estudantes, pois também é divulgada por todos os meios de comunicação.
Trata -se de tema geral que, no entanto, normalmente, não se inclui nas grades curriculares de cursos superiores que não sejam da área de Saúde.
Apesar disso, a temática é de conhecimento geral e relacionada aos direitos humanos, esses sim, ponto de integração com as diretrizes curriculares.
A linguagem utilizada era clara, indicando a existência dos fatos vinculados a questões relevantes e à saúde pública.
No entanto, como os textos trazem uma ampla gama de possibilidades de interpretação, ocorreram algumas compreensões diversas das esperadas.
Além disso, ressalta -se que a amplitude dos textos disponibilizados no enunciado fornece elementos que poderiam ser simplesmente reproduzidos pelos concluintes nas suas respostas.
O enunciado da questão encaminha a resposta para dois aspectos.
O estudante deveria abordar conteúdo pertinente à vulnerabilidade das mulheres às DST e o papel social do homem em relação à prevenção dessas doenças.
Para responder a esses aspectos, um número expressivo de concluintes utilizou somente os elementos dos textos motivadores, sem nada acrescentar.
A resposta deveria conter também uma segunda parte, a indicação de duas ações voltadas para o público masculino, a serem adotadas no âmbito das políticas públicas de saúde ou de educação, para a redução do problema.
Verificou -se que o desempenho dos estudantes nessa questão foi de médio para bom.
Ao se analisarem as respostas, constata -se que muitos foram os que apresentaram os tópicos presentes no padrão de resposta, amplo o suficiente para albergar as respostas apresentadas pelos concluintes.
Tal resultado pode ser justificado, tendo em vista tratar -se de tema bastante difundido nas diversas mídias, por ter sido apresentado de forma ampla pelos textos motivadores presentes no enunciado.
Por outro lado, o debate no ambiente universitário sobre as relações de poder entre homens e mulheres tende a ser marcante pelo movimento de empoderamento feminino, liderado por organizações sociais com ampla atuação nos estabelecimentos de nível superior.
De forma recorrente, as relações patriarcais entre homens e mulheres foram vinculadas ao modelo machista vigente.
Essas condições, segundo os concluintes, explicam a incidência da doença e a situação de vulnerabilidade das mulheres, bem como a função social dos homens.
Outras respostas apontaram o poder público, nas áreas de saúde e/ou educação, como órgãos responsáveis pela solução do problema.
Para os que se fixaram na área da saúde, foram frequentes a indicação da necessidade de mais medicamentos, vacinas, educação sexual, médicos e postos de atendimentos.
Alguns, no entanto, apenas indicaram ações individuais para o homem, e outros, ações individuais para as mulheres.
Também figuraram respostas que apresentavam ações da igreja como solução, além de outras que faziam referências a campanhas no seio das famílias. 6.3.1.3 Análise de Conteúdo da Questão Discursiva 2 do Componente de Formação Geral A Tabela 6.15 apresenta as estatísticas básicas das notas dos alunos na questão discursiva 2 da área de Ciência da Computação por Grande Região, segundo Habilitação e Modalidade de ensino.
Como já comentado, o desempenho médio dos estudantes de cursos de Ciência da Computação, tanto de Licenciatura, quanto de Bacharelado, na questão discursiva 2 foi ligeiramente inferior ao obtido na questão discursiva 1 : médias, respectivamente 52,6 e 52,1 na questão 1 e médias 50,4 e 52,0, na questão 2.
Entre os cursos de Licenciatura, a região Sul foi aquela cuja Média nessa questão foi maior ( 61,7 ), e a de menor Média foi a região Centro-Oeste ( 43,6 ).
Quando se considera também a Modalidade de ensino, a Tabela 6.15 mostra que o desempenho médio dos estudantes de cursos Presenciais de Licenciatura na questão discursiva 2 ( média 50,8 ) foi inferior ao obtido na questão discursiva 1 ( média 52,4 ).
A região Sul foi aquela cuja Média nessa questão foi maior ( 61,5 ), e a de menor Média foi a região Centro-Oeste ( 45,3 ).
Quanto à variabilidade das notas, o Desvio padrão de todo o Brasil foi 33,8.
O maior desvio nessa questão foi obtido na região Nordeste ( 34,9 ), enquanto o menor foi obtido na região Sul ( 29,8 ).
A Mediana para esses alunos de cursos Presenciais de Licenciatura foi 50,0, a mesma para todos as regiões, com exceção da região Sul ( 70,0 ).
As notas máximas ( 100,0 ) e as notas mínimas ( zero ) foram as mesmas em todas as regiões.
A Média das notas dos estudantes de cursos a Distância de Licenciatura na questão discursiva 2 foi 49,5.
A região Sul foi aquela cuja Média nessa questão foi maior ( 72,5 ), e a de menor Média foi a região Centro-Oeste ( 37,0 ).
O Desvio padrão de todo o Brasil foi 34,1, sendo que o maior desvio nessa questão foi obtido na região Sul ( 38,9 ), enquanto o menor foi obtido na região Norte ( 28,3 ).
A Mediana para estes alunos de cursos a Distância de Licenciatura foi 50,0, a mesma das regiões Nordeste e Sudeste.
As notas máximas ( 100,0 ) e as notas mínimas ( 0,0 ) foram as mesmas em todas as regiões, com exceção da Máxima da região Centro-Oeste ( 75,0 ) e da Mínima da região Sul ( 45,0 ).
A desagregação por Grande Região para as notas de Bacharelado, na questão discursiva 2, mostra que a maior Média foi observada na região Sudeste ( 53,3 ), e a menor, na região Centro-Oeste ( 47,5 ).
O Desvio padrão das notas do conjunto desses alunos foi 33,9, sendo que o maior desvio nessa questão foi obtido na região Norte ( 35,6 ), enquanto o menor foi obtido na região Sudeste ( 33,1 ).
A Mediana para esses alunos de Bacharelado foi 50,0, a mesma em duas regiões : Norte e Centro-Oeste.
As notas máximas ( 100,0 ) e as notas mínimas ( zero ) foram as mesmas em todas as regiões.
Como já comentado, não há registro em cursos a Distância de Bacharelado.
Todos os registros coincidem no Total e na Educação Presencial.
O Gráfico 6.8 mostra a distribuição das notas na questão discursiva 2 do Componente de Formação Geral.
Observa -se que a maior frequência corresponde aos alunos que obtiveram nota dentro do intervalo ( 70 ; 80 ] para ambas as Habilitações.
Destacam -se, também, os alunos que deixaram a resposta em branco, com proporção entre 15 \% e 20 \% para ambas as Habilitações. 6.3.1.4 Comentários sobre a correção de Conteúdo das respostas à Questão Discursiva 2 O enunciado era claro e abordava um tema de extrema relevância : a questão de gênero.
Trata -se de uma proposta polêmica, pois abrange a análise da inserção social de minoria vinculada à questão de gênero.
O enunciado solicitava que o concluinte discorresse sobre a importância do nome para as pessoas transgêneras e propusesse uma medida, no âmbito das políticas públicas, para facilitar o acesso dessas pessoas à cidadania.
A respeito da primeira parte, a importância do nome, esperava -se que o estudante abordasse as dificuldades que o uso do nome civil ( nome de registro, nome de batismo, nome da certidão de nascimento, por exemplo ) pode acarretar para as pessoas transgêneras, além de demonstrar conhecimento sobre a importância do nome social, direito já reconhecido através, principalmente, de decisões judiciais e atos administrativos, como decretos e portarias que dispõem sobre o uso do nome social.
Em atendimento à segunda parte do comando, o participante deveria propor uma política pública associada ao direito à cidadania das pessoas transgêneras.
Sobre esse aspecto, esperava -se a compreensão do significado de política pública, ou seja, propostas de âmbito governamental, e não, iniciativas de Organizações Sociais, igrejas ou grupos sociais de qualquer tipo.
O nascimento das pessoas, no Brasil, é objeto de inscrição em cartórios de Registro Civil das Pessoas Naturais, órgão responsável por estabelecer a certificação pública dos dados de todos os indivíduos que não somente nascem, mas casam, eventualmente se divorciam e vêm a óbito.
Daí o nome civil.
Para flexibilizar a correção, outras nomenclaturas foram consideradas, a saber : identificação de nascimento, nome de batismo, nome de registro, nome próprio, nome verdadeiro, que diferenciam o nome civil daquele que o indivíduo deseja utilizar para traduzir a sua opção pessoal de gênero, que se convencionou denominar nome social.
Quanto ao entendimento do termo transgênero, deve -se destacar que o movimento de mudança do gênero é contemporâneo e não está necessariamente vinculado à mudança de sexo.
Existem várias hipóteses em jogo.
Uma é a que envolve cirurgias de mudança de sexo ( a menina extirpa os seus seios e inclui órgão genital por intervenção cirúrgica ; o menino inclui seios e extirpa o seu órgão genital e inclui o órgão feminino também por intervenção cirúrgica ).
Há outro movimento que envolve apenas as vestimentas e o linguajar corporal, além dos relacionamentos ( as meninas assumem o trajar, o andar, os gostos dos meninos e mantêm relacionamentos afetivos com meninas ; os meninos assumem o trajar, o andar, os gostos das meninas e mantêm relacionamentos afetivos com meninos ).
Em relação à discussão da primeira solicitação do comando da questão, destaca -se que todos têm nome civil, que foi escolhido pelos seus pais ou responsáveis e, na falta deles, por autorização do Juiz de Direito, pautado no sexo detectado no nascimento.
Já o nome social é aquele que indica a preferência do indivíduo na sua identificação nas relações sociais e que destoa do seu gênero de registro civil.
Ressalte -se que o Brasil não possui ainda legislação específica sobre pessoas transgêneras.
Observa -se um relativo progresso no reconhecimento de direitos, através, principalmente, de decisões judiciais e atos administrativos, como decretos e portarias que dispõem sobre o uso do nome social.
Sem uma lei que defina os procedimentos da alteração dos documentos para pessoas transexuais, essa parcela da população LGBT ( lésbicas, gays, bissexuais, transexuais e transgêneros ) é obrigada a procurar na Justiça o reconhecimento de sua identidade, em processos que podem ser longos e que dependem do olhar de determinados, profissionais, como médicos, psicólogos, assistentes sociais, bem como do entendimento dos advogados, defensores públicos, juízes e do Judiciário, como um todo.
Atualmente, tramita na Comissão de Direitos Humanos da Câmara dos Deputados o Projeto de Lei João Nery ( 5002/2013 ), dos deputados Jean Willys ( PSOL-RJ ) e Erika Kokay ( PT-DF ), que determina que o reconhecimento da identidade de gênero seja um direito do cidadão.
O projeto recebeu o nome do primeiro trans-homem operado no Brasil.
Grande parte dos respondentes considerou a incongruência entre a imagem corporal, representada pelo gênero assumido pelos trans, e o nome que apresenta nos seus documentos como causa de sérios constrangimentos e embaraços e as dificuldades para que esses sujeitos acessem direitos básicos, como saúde, educação e empregabilidade em uma sociedade heteronormativa e preconceituosa.
Muitos demonstraram reconhecer que a identidade e expressão de gênero, através do nome social, pressupõe que sejam respeitadas a dignidade da pessoa humana, a intimidade, a liberdade e o direito à cidadania plena.
No entanto, também foram encontradas, em quantidades significativas, respostas que demonstravam a existência de uma interpretação social conservadora que não admite a mudança de sexo nem do nome social.
Existe um acentuado preconceito contra todos os movimentos que buscam identificação social diversa da civil.
Isso inclui os variados grupos : gays, travestis, transgêneros, dentre outros.
Já em relação ao segundo tópico a ser abordado pelo concluinte, é notório que as pessoas trans encontram -se destituídas de direitos no campo das políticas públicas, como saúde, habitação, emprego, educação, justiça, dentre outros.
Nesse sentido, uma grande diversidade de medidas, no âmbito das políticas públicas, poderia ser proposta.
Ações voltadas para a capacitação de profissionais da educação, saúde, segurança pública, dentre outras áreas, para melhor atender às demandas e especificidades desse público alvo seriam necessárias.
Ainda em relação à educação, seriam também necessárias ações que contribuíssem para diminuir a evasão, pois o ambiente escolar é difícil para uma pessoa transgênera, por conta das situações discriminatórias e preconceituosas vindas dos colegas de classe, e mesmo das/dos professores e/ou de outros profissionais da área, não podem ser mais um complicador nesse processo.
Ainda no que tange ao assunto, percebe -se que, ao serem excluídas dos bancos escolares, as oportunidades no mercado de trabalho formal praticamente inexistem.
Um outro aspecto que se evidencia no cotidiano desse segmento da sociedade diz respeito às situações de violência sofridas pelas pessoas transgêneras.
Nesse sentido, algumas propostas foram elencadas para minimizar o problema e combater a violência.
Políticas envolvendo melhor formação de funcionários das polícias civil e militar, por exemplo, somariam esforços na qualificação das abordagens policiais e dos atendimentos prestados pelas forças de segurança, pois o respeito à orientação sexual e à identidade e expressão de gênero é, antes de tudo, o respeito à dignidade, à cidadania e à própria democracia.
Nesse sentido, revela -se essencial que, ao sofrerem agressões, sejam elas físicas, psicológicas ou morais no espaço público ou domiciliar, as pessoas trans possam ser atendidas, com respeito inclusive ao nome social, nas delegacias comuns e, inclusive, nas Delegacias Especiais de Atendimento à Mulher ( DEAM ), pois esse é o órgão institucional que possivelmente detém as qualificações necessárias para o atendimento ao gênero feminino ao qual as travestis e mulheres trans efetivamente pertencem.
Também ocorre uma latente inacessibilidade das pessoas trans à assistência jurídica.
Nesse sentido, é necessário implementar projetos de capacitação para profissionais do poder judiciário, como um todo, de modo a prestarem a assistência jurídica gratuita na defesa intransigente dos direitos sociais, civis, trabalhistas e políticos dessa população.
Dessa maneira, a Justiça agiria de forma mais acessível e sensível para o segmento, que terá suas especificidades observadas a começar pelo devido respeito do direito ao uso do nome social.
Nesse âmbito, o mais comum foi propor a adoção de medidas punitivas para quem viola o direito à autodeterminação de gênero e a elaboração de leis que garantam a mudança do nome e assegurem outros direitos para as pessoas trans.
Propostas envolvendo a formação de profissionais, como as mencionadas aqui, e de conscientização da população foram as mais frequentes.
Muitas outras enfatizavam o direito ao uso respeitado do nome social e a garantia de alteração do registro civil, independentemente da cirurgia de redesignação sexual ou transgenitalização.
Muitos questionaram o fato de não bastar a autodeclaração no processo de alteração de nome, enfatizando que o nome social é uma questão de princípio de dignidade e cidadania.
Apesar de pouco mencionada nas respostas corrigidas, políticas envolvendo a promoção da saúde são muito necessárias.
Cabe aos gestores de saúde ações direcionadas, não somente à prevenção e tratamento do HIV/Aids, mas também políticas de saúde que contemplem as especificidades dessa população, tais como o Ministério da Saúde, especialmente o SUS ( Sistema Único de Saúde ), incluir, de fato, o atendimento para pessoas trans para acompanhamento da terapia hormonal, que muitas demandam, respeitando -se o direito de utilização do nome social.
De modo geral, são muitas as políticas públicas a que precisam ser implantadas.
Urge a necessidade de uma agenda mais compromissada por parte dos gestores do nosso país com essa comunidade, que ainda é muito discriminada.
Logo, tais políticas devem contemplar suas necessidades básicas como o direito de acesso aos estudos, à capacitação para o mercado de trabalho, à profissionalização, ao acesso a bens e serviços de qualidade em saúde, habitação, segurança, cultura e assistência social e, certamente, o respeito ao uso do nome social, o que também contribui para essa inclusão.
Assim, as opções apresentadas no padrão de resposta funcionaram como exemplos de políticas públicas possíveis, já que seria difícil apresentar uma lista suficientemente ampla para abranger a grande variedade de respostas encontradas.
O desempenho dos estudantes nessa questão foi médio.
As respostas, de modo geral, indicaram a necessidade de uma atualização por parte dos concluintes quanto a temas essenciais e, no caso da presente questão, a análise sobre aspectos fundamentais da questão de gênero.
Muitos, possivelmente motivados pelos textos apresentados no enunciado, restringiram -se a discorrer sobre preconceito de gênero ou mesmo sobre o preconceito de modo mais geral, sem responder ao que foi solicitado no comando da questão.
Outros apresentaram discursos religiosos ou ideológicos, também sem tratarem do que foi solicitado. 6.3.1.5 Análise de Língua Portuguesa das Questões Discursivas do Componente de Formação Geral Os dados obtidos a partir das respostas às questões discursivas do Componente de Formação Geral de Ciência da Computação desagregados por Grande Região, segundo Habilitação e Modalidade de Ensino, no que tange à Língua Portuguesa, encontram -se na Tabela 6.16 e no Gráfico 6.9.
Nos cursos de Licenciatura em Ciência da Computação, a Média desse aspecto foi 53,4, ligeiramente acima da obtida pelos alunos de Bacharelado ( 53,0 ).
Quando se considera a Modalidade de ensino para cada Habilitação, os alunos de cursos Presenciais de Licenciatura obtiveram Média 53,3, um pouco menor do que a Média com respeito à Língua Portuguesa para os alunos de cursos a Distância ( 53,4 ).
Para os cursos de Bacharelado, a Média foi 53,0, como já comentado.
A Mediana das notas de Língua Portuguesa para os cursos de Licenciatura foi 60,0 e também para ambas modalidades de ensino.
Para os cursos de Bacharelado a Mediana foi 60,0.
A nota Máxima não chegou a 100,0 para nenhuma combinação de Habilitação e Modalidade.
O Gráfico 6.9 mostra a distribuição das notas de Língua Portuguesa do Componente de Formação Geral.
Observa -se que a maior frequência ( pouco menos do que 30 \% ) corresponde aos alunos que obtiveram nota no intervalo ( 60 ; 70 ], nas duas Habilitações.
Destacam -se, também, os alunos que deixaram as respostas a ambas as questões em branco, representando um pouco mais do que 10 \% do total, caracterizando -se como máximo local. 6.3.1.6 Comentários sobre a correção das respostas de Formação Geral com respeito à Língua Portuguesa Os enunciados apresentados em 2017 não explicitaram a exigência de elaboração de um texto “ dissertativo ”, como nos anos anteriores.
O participante deveria fazer uma exposição de seus conhecimentos sobre cada um dos assuntos e estruturar seus textos de acordo com as características do registro formal adequado à situação comunicativa – avaliação de conhecimentos.
Essa configuração determina exigências quanto aos seguintes aspectos : adequação da seleção vocabular, desenvolvimento do conteúdo, organização lógica das ideias, estruturação sintática dos períodos, utilização de procedimentos de encadeamento textual e de referenciação, obediência às exigências morfossintáticas próprias da modalidade escrita da norma-padrão, respeito às regras ortográficas e às regras de acentuação gráfica.
O padrão de resposta utilizado na avaliação das questões 1 e 2 considerou os aspectos relevantes ao bom desempenho linguístico como competências distintas, de modo a permitir um mapeamento detalhado do domínio dos recursos disponíveis na Língua Portuguesa para a comunicação escrita formal.
Com base nesse objetivo, foram avaliados os seguintes aspectos : ( a ) estruturação textual condizente com o gênero solicitado e o modo de organização textual expositivo adequado ao gênero – essa competência envolve : • estruturação sintática condizente com o padrão da modalidade escrita formal da língua portuguesa de modo a garantir a clareza necessária ; • distribuição do conteúdo do texto em parágrafos, de modo a garantir a sua organização temática ; • utilização de operadores discursivos que contribuam para a progressão temática do texto, estabelecendo relações lógicas entre as ideias apresentadas, tanto do ponto de vista intrafrasal, como do interfrasal ; • utilização de procedimentos de referenciação lexical e pronominal que permitam a retomada de referentes textuais ; • utilização de sinais de pontuação que contribuam para a organização lógica da frase e do texto ; • inteligibilidade relacionada ao atendimento das exigências de estruturação textual.
Espera -se, portanto, que o participante recorra a procedimentos linguístico-discursivos para organizar seu texto, permitindo o encadeamento lógico entre suas partes de forma a garantir a progressão e a coerência textuais.
Isso significa que os seguintes procedimentos foram considerados inadequados, de acordo com o padrão de resposta proposto : • elaboração de frases fragmentadas que comprometam a estrutura lógico-gramatical do texto ; • sequência justaposta de ideias sem encaixamentos sintáticos, reproduzindo hábitos da oralidade ; • elaboração de frase com apenas oração subordinada, sem oração principal ; • emprego equivocado do conector ( preposição, conjunção, pronome relativo, alguns advérbios e locuções adverbiais ) comprometendo a expressão da relação lógica entre duas ideias, com prejuízo da clareza do texto ; • emprego do pronome relativo sem a preposição, quando obrigatória ; • repetição ou substituição inadequada de palavras sem o emprego dos recursos oferecidos pela língua ( pronome, advérbio, artigo, sinônimo ) ; • emprego inadequado dos pronomes relativos “ cujo(a) ” e “ onde ” ; • utilização inadequada dos sinais de pontuação, comprometendo a clareza textual ; • inteligibilidade relacionada ao atendimento das exigências de estruturação textual. ( b ) respeito às convenções ortográficas da norma-padrão da Língua Portuguesa – essa competência envolve o domínio das regras de acentuação gráfica e da grafia padrão das palavras ( com ausência de abreviaturas próprias da linguagem da internet ), de acordo com as convenções estabelecidas pela legislação em vigor e consubstanciadas no Vocabulário Ortográfico da Língua Portuguesa, editado pela Academia Brasileira de Letras. ( Vale notar que, nessa edição de 2017, não foi aceita a legislação anterior, no caso das regras relativas ao uso do hífen e da acentuação gráfica ).
Espera -se que o participante : • grafe corretamente as palavras ; • respeite as regras de acentuação gráfica ; • empregue maiúsculas em início de frase, em nomes próprios de pessoas, lugares ou instituições ; • grafe as siglas com letras maiúsculas ou apenas com a primeira letra maiúscula quando formarem uma palavra ; • evite abreviações como p/, vc, tb, pq, tá, né, usadas muitas vezes em escrita informal e na internet ; • obedeça às regras de separação de sílabas no final da linha. ( c ) domínio dos diferentes aspectos morfossintáticos próprios da modalidade escrita formal da norma-padrão da Língua Portuguesa – essa competência envolve : a concordância nominal, a concordância verbal, a regência nominal, a regência verbal, a flexão nominal, a flexão verbal, a correlação entre os tempos verbais, a colocação pronominal e a utilização de sinais de pontuação que contribuam para a organização lógica da frase e do texto.
Espera -se que o participante : • flexione o verbo para estabelecer concordância de número com o sujeito da frase ; • flexione o artigo, o adjetivo e o pronome para concordar em número e em gênero com o substantivo a que se referem ; • observe a regência nominal e a verbal, utilizando a preposição adequada depois de um substantivo, um verbo ou um adjetivo ; • empregue adequadamente o acento grave indicador da crase entre uma preposição e um artigo ( a+a ) ; • obedeça às regras de colocação pronominal ( próclise e ênclise ), distintas dos hábitos da oralidade ou da escrita informal ( exigência de próclise com termo atrator, não exigência de que o pronome oblíquo se ligue ao verbo auxiliar por meio de hífen ) ; • flexione adequadamente verbos, substantivos, adjetivos e pronomes no que diz respeito à expressão das categorias gramaticais ;  flexione os verbos para expressar a correlação de modo e tempo nas estruturas subordinadas.
Com base nesses critérios, foram considerados como desvios de caráter morfossintático, e não como desvios ortográficos, as alterações que envolvem mudança de classe gramatical ou de forma flexional do verbo :  eliminação da marca de infinitivo ( -r- ) e substituição por acento agudo ou ausência total de marca do infinitivo ;  confusão entre “ ão ” e “ am ” nas formas verbais ;  confusão entre “ há ” e “ a ” ;  uso de hífen para separar pronome átono – tanto uso indevido quanto omissão ( exemplo : “ esperasse ”, em lugar de “ espera -se ” ; “ falar -mos ”, no lugar de “ falarmos ” ;  “ esta ” ( no lugar de “ está ” ) ; “ mais ” ( no lugar de “ mas ” ) ; “ e ” ( no lugar de “ é ” ).  verbos “ ter ” e “ vir ” que, na terceira pessoa do plural, não apresentarem o acento circunflexo, serão considerados como desvio de concordância, em aspectos morfossintáticos. ( d ) seleção vocabular adequada à modalidade escrita formal da Língua Portuguesa, exigida pela situação comunicativa – essa competência envolve a precisão na seleção/utilização do vocabulário relacionado à temática solicitada pela questão ; a ausência de marcas da oralidade, como termos de sentido muito genérico ( “ coisa ”, “ negócio ”, “ você ” ) e termos de registros mais informais ( como gírias, jargões, frases feitas, ditados populares, termos regionais ).
Assim, espera -se que o participante respeite a adequação vocabular não usando gírias ou expressões coloquiais, evite repetição desnecessária de palavras e utilize um vocabulário mais formal, como solicitado por um texto dissertativo.
Observações : - A inteligibilidade relacionada ao atendimento das exigências de estruturação textual foi avaliada na segunda competência ( relativa aos aspectos textuais ). - Os problemas de coerência textual provocados por uso indevido do vocabulário foram avaliados na terceira competência ( relativa aos aspectos morfossintáticos e vocabulares ). - Não foram considerados como desvios morfossintáticos os problemas de caligrafia ( -a/ -o, -s, -r ). - Não foram considerados os textos grafados integralmente em caixa alta, já que não seria possível distinguir alguma marcação especial para as letras em início de frase. - Cada desvio foi considerado como uma ocorrência, mesmo que dois desvios fossem relativos ao mesmo aspecto linguístico.
A grade de avaliação do desempenho linguístico considerou, portanto, três grandes grupos de competências, segundo os aspectos explicitados anteriormente : a ) domínio das convenções ortográficas : grafia de vogais e consoantes, uso de maiúsculas e minúsculas, emprego do hífen e acentuação gráfica ; b ) domínio dos procedimentos de estruturação textual do ponto de vista microestrutural : organização interna dos períodos, emprego de conectores para a articulação lógica entre os períodos e entre os parágrafos, emprego de marcas de referenciação lexical e pronominal ; utilização dos sinais de pontuação que contribuem para a organização lógica da frase ; c ) domínio das regras de caráter morfossintático estabelecidas como modelares do ponto de vista da modalidade escrita formal da norma-padrão da Língua Portuguesa : concordância nominal e verbal, regência nominal e verbal, colocação pronominal, flexão nominal e verbal, correlação entre tempos e modos verbais, ausência de marcas de oralidade.
A seleção vocabular adequada à modalidade escrita formal da Língua Portuguesa foi incorporada a essa última competência, tendo em vista a intersecção entre as duas do ponto de vista das exigências do registro formal da modalidade escrita da norma-padrão.
Apreciação geral do desempenho dos estudantes : A correção revelou desempenhos distintos dos participantes, marcados pela falta de repertório cultural da maioria dos participantes : - a questão 1 propiciou a oportunidade de maior desenvolvimento do tema solicitado, gerando consequências na estruturação textual, já que os textos foram mais longos, mais elaborados e fluentes.
Entretanto, em virtude da diversidade de abordagens dos textos motivadores ( a falta de penicilina, a recusa no uso da camisinha pelos homens, a maior frequência das mulheres na realização de exames ginecológicos, a transmissão das DST aos bebês durante a gestação ), observaram -se respostas que se configuraram como verdadeiras paráfrases, com comprometimento do caráter autoral desejável em uma questão discursiva de Formação Geral.
Alguns participantes selecionaram um trecho de cada texto motivador para compor suas respostas, outros copiaram trechos inteiros. - a questão 2 teve um comportamento distinto, devido ao caráter polêmico que o caracteriza.
Assim, os textos são mais curtos, com vocabulário repetitivo e preso aos textos motivadores.
Além disso, observou -se grande ocorrência de protesto em função do tema, ora por motivos religiosos ora por preconceito de gênero.
Quanto aos aspectos linguísticos analisados durante esta avaliação, os resultados observados estão descritos nos parágrafos seguintes.
Aspectos ortográficos : O desempenho dos participantes revelou uma diferença muito grande nos dois aspectos analisados nesta competência : baixo índice de desvios da grafia padrão e grande índice de desvios de acentuação.
Em vários casos, ocorre ausência completa de acentuação gráfica.
Os resultados revelam que a tendência dominante entre os universitários brasileiros é a eliminação da acentuação gráfica, talvez motivada pelos hábitos relacionados às redes sociais e pela ausência de esclarecimento dos meios de comunicação, das autoridades e das escolas sobre as decisões do Acordo Ortográfico da Língua Portuguesa de 1990.
Os casos mais sistemáticos de eliminação do acento indicador da sílaba tônica são : • palavras proparoxítonas ( “ sifilis ”, “ proximos ”, “ politicas ”, “ publicas ”, “ transgeneros ” ) ; • palavras paroxítonas terminadas em ditongo crescente ( “ necessario ”, “ noticiarios ”, “ individuo ”, “ dependencia ”, “ varios ”, “ propria ”, “ transmissiveis ”, “ ocorrencia ” ) ; • palavras paroxítonas com hiato ( “ saude ” ) • palavras oxítonas ( “ ninguem ”, “ esta ”, “ ate ”, “ tambem ”, “ prevencao ” ).
Por outro lado, destaca -se o uso indevido do acento gráfico em palavras como “ genêro ”, “ melâncolia ”, “ prevênção ”, “ intervênção ”, por exemplo.
Há, também, uma tendência a acentuar paroxítonos como se fossem oxítonos : “ transmissivéis ”.
Quanto ao domínio das convenções relativas à grafia das palavras, observam -se desvios recorrentes, como a hipercorreção pela escolha de “ e ” no lugar de “ i ”, por influência de hábitos da oralidade ( “ descriminação ” no lugar de “ discriminação ”, “ entervenção ” no lugar de “ intervenção ” ) ; desvios de grafia relacionados à variação diastrática podem ser observados em “ estrupo ”, “ subjulgadas ”, “ precoseito ”, “ soubre ”, “ vecendo ”, “ indesencia ”, “ apolojia ”, “ fulga ”, “ dereitos ”, “ sifelis ”.
Observam -se, também, casos de inadequação no uso da maiúscula para destacar determinadas palavras-chave do texto, como “ Transgêneros ”, “ Brasileiros ”, “ Homens ”.
Destaque -se, também, o grande número de participantes que grafam os textos inteiramente em caixa alta.
Vale observar, também, que, ao contrário do que se esperava, não apareceram abreviaturas próprias do “ internetês ”, relacionadas ao uso de redes sociais e e-mails.
Aspectos textuais : Esta competência é a que se revela como a mais problemática entre os participantes, porque são muitos os problemas observados, desvios acumulados durante toda a formação do estudante e que não se resolvem com um estudo autodidata, como acontece com regras ortográficas ou morfossintáticas.
São eles : sequência justaposta de ideias sem encaixamentos sintáticos ; redução drástica de estruturas subordinadas, ao lado do aumento na frequência de estruturas coordenadas e absolutas ; redução no uso de conectores para expressar relações lógicas essenciais à construção do texto, substituídas pela exigência de inferência por parte do interlocutor para suprir a sua ausência ; emprego equivocado de operadores que não estabelecem relações lógicas coerentes entre ideias do texto ; emprego inadequado do pronome relativo ( com omissão da preposição ou a utilização de pronome inadequado, como “ onde ” ) ; repetição exaustiva de termos sem a utilização de procedimentos mais sofisticados de substituição ( hiperonímias, hiponímias, nominalizações, expressões metafóricas ) ; frases fragmentadas que comprometem a estrutura lógico-gramatical ; frases formadas apenas por oração subordinada, sem oração principal.
Um importante aspecto a destacar é o baixo desempenho de uma parte dos participantes em relação à estrutura formal do texto produzido, o que é preocupante ao se levar em conta que são graduandos em fase final de formação.
São frequentes os casos de desvios de estruturação frasal, com uso inadequado ou ausência de conectivos entre parágrafos e entre frases.
Em uma parte dos textos, falta textualidade e domínio do registro padrão da língua.
Na verdade, observam -se relações linguísticas quase agramaticais, como as estabelecidas pela sequência de gerúndios sem o apoio de um ponto de partida para a organização das informações gramaticais e semânticas.
Observou -se que uma parte dos participantes não distribuiu as ideias em parágrafos, talvez devido ao pequeno número de linhas disponibilizadas para a resposta da questão ou, quem sabe, pela suposição de que não seria necessária essa divisão por não se tratar de um texto no modelo de uma redação dissertativo-argumentativa, como solicitado nos vestibulares.
Em função do encaminhamento dos enunciados das duas questões, que solicitaram uma análise do problema e encaminhamentos de políticas públicas, houve uma grande tendência, também, de construção de dois parágrafos desconexos, sem utilização de elementos coesivos adequados à progressão textual.
Quanto à utilização dos mecanismos de referenciação, deve -se destacar a ocorrência de repetições de palavras ou expressões sem a utilização de termos sinônimos ou pronomes, como seria adequado.
Quanto à utilização dos sinais de pontuação, observou -se uma grande precariedade nos textos analisados.
É muito frequente a ocorrência de parágrafos sem marca interna de pontuação para separar os períodos.
Vale observar que não foi penalizada a ausência de vírgula para destacar locuções ou adjuntos adverbiais de pequena extensão deslocados de posição na frase, por ser um uso opcional.
São os seguintes os tipos de problemas encontrados : a ) vírgula : utilização de vírgula para separar o sujeito e o predicado ; ocorrência de apenas uma das vírgulas para separar uma palavra, uma expressão ou uma oração encaixada ; uso de vírgula no lugar do ponto para separar ideias que constituem períodos distintos ; ausência de vírgula para separar elementos de uma enumeração ; ausência de vírgula para separar oração adjetiva explicativa ou utilização inadequada para separar oração adjetiva restritiva ; b ) ponto e vírgula : utilização do ponto e vírgula no lugar de vírgula ; c ) ponto final : ausência de ponto final para separar períodos.
Aspectos morfossintáticos e vocabulares : Em relação à regência, o desvio mais frequente é a falta do sinal indicativo da crase – isso revela que o usuário não tem consciência de que, sob a forma do termo “ a ”, existe a presença de uma contração entre a preposição “ a ” ( exigida pela regência do termo anterior ) e o artigo definido “ a ”.
Um desvio de regência significativo, nos últimos anos, é a utilização inadequada de uma preposição ou sua ausência após o verbo ou o nome ( substantivo ou adjetivo ).
Outro problema relacionado à regência verbal e à nominal, encontrado frequentemente nas questões, foi a ausência de preposição antes de pronome relativo, processo generalizado na modalidade oral da língua, em situações de registro informal.
Apesar da possibilidade de que essa alteração de regência se generalize no padrão escrito da Língua Portuguesa, como já está ocorrendo até em textos jornalísticos, o não emprego da preposição foi considerado inadequado neste processo de avaliação.
A concordância verbal e a concordância nominal apresentam alguns desvios muito frequentes.
Quanto à concordância de número, observou -se ausência de marca ( com sujeito anteposto ou posposto ) ou uso indevido ( uso inadequado da marca de plural comandado pelo núcleo plural da locução adjetiva, apesar de o substantivo que funciona como núcleo do sintagma nominal estar no singular ).
Uma ocorrência que se destacou foi a ausência de acento circunflexo na forma plural do presente do indicativo do verbo “ ter ”, que foi considerada como um desvio na concordância verbal e não na acentuação gráfica.
Quanto à concordância de gênero, vários casos foram observados, normalmente no âmbito de sintagmas nominais longos, em que o adjetivo está afastado do substantivo.
Deve -se destacar o aparecimento da marca de plural em verbos ou adjetivos comandados por adjunto adnominal plural, apesar de serem relacionados a núcleos substantivos no singular, evidenciando um processo de hipercorreção ( exemplo : “ O nome social dos trangêneros podem ser estabelecidos por uma legislação específica ” ).
Quanto à questão da colocação pronominal, foram poucos os desvios observados.
Concluiu -se que, no registro escrito formal, a maioria dos participantes já incorporou regras como a não introdução da frase por um pronome oblíquo e a próclise na presença de um termo atrator.
Não se adotou, entretanto, o padrão excessivamente formal descrito pelas gramáticas normativas em relação à posição do pronome oblíquo em locuções verbais, já que esse uso está muito distante da prática cotidiana, até em textos mais formais.
Quanto aos aspectos vocabulares, alguns tipos de inadequação foram observados : expressões da oralidade ; seleção vocabular incompatível com o contexto, gerando falta de inteligibilidade ; falta de domínio de vocabulário mais abstrato e de maior complexidade, essencial ao desenvolvimento do texto de base dissertativa.
O principal aspecto observado foi a excessiva repetição de certas palavras, revelando limitação de repertório vocabular.
O termo “ pessoa ”, por exemplo, chega a ser repetido até 7 ou 8 vezes em um mesmo texto.
Várias marcas de oralidade foram identificadas, embora não com alta frequência : uso do pronome relativo “ onde ” como relativo universal, falta de artigo definido antes de substantivo, reduções como “ tá ”, “ pra ”, “ pro ”, “ prum ”, expressões informais, eliminação de preposições.
Em função do tema solicitado na questão 1, a sigla DST foi grafada de diferentes maneiras : DST´s, DSTs ou DSTS.
Além disso, houve flutuação de gênero gramatical nos sintagmas que envolviam a sigla : “ os diferentes DSTs ” ou “ as diferentes DSTs ”.
Em função do tema solicitado na questão 2, houve flutuação de gênero gramatical nos sintagmas que envolviam a sua palavra-chave : “ os trans ” ou “ as trans ” ; “ os transgêneros ” ou “ as transgêneras ”.
Todas essas formas foram aceitas, por considerarmos que sua utilização é muito recente e os participantes ainda não sistematizaram essas formas em seu uso da língua. 6.3.2 Componente de Conhecimento Específico As estatísticas básicas das notas referentes às questões discursivas do Componente de Conhecimento Específico de Ciência da Computação Licenciatura, por Grande Região, segundo Modalidade de Ensino, estão disponibilizadas na Tabela 6.17a.
As informações equivalentes para o Bacharelado estão na Tabela 6.17b.
Observa -se que as Médias foram bem mais baixas do que para as questões discursivas do Componente de Formação Geral.
Enquanto no Componente de Formação Geral as Médias para estudantes de Ciência da Computação de Licenciatura foram, respectivamente para os de Ensino Presencial e a Distância, 51,9 e 51,5 ( 51,8 para o total ), na parte de Conhecimento Específico as Médias foram 17,1 e 14,7 ( 16,3 para o total ).
Já para os cursos de Bacharelado, a Média foi 52,2 ( não houve registro de Educação a Distância no Bacharelado ) para as discursivas de Formação Geral e 28,5 para as de Conhecimento Específico.
Para os cursos de Licenciatura, a nota Máxima foi maior para os cursos Presenciais ( 90,0 ) do que para os a Distância ( 65,0 ).
Para os cursos de Bacharelado a nota Máxima foi 100,0.
As notas mínimas foram zero para todas as combinações de Habilitação e Modalidade de Ensino, com exceção da Mínima da região Sul da Educação a Distância da Licenciatura que foi 16,7.
O Gráfico 6.10 representa a distribuição das notas nas questões discursivas no Componente de Conhecimento Específico.
A moda desta distribuição ocorre no intervalo [ 0 ; 10 ], com cerca de 40 \% do total de participantes da modalidade Bacharelado e mais de 45 \% da habilitação Habilitação em Licenciatura.
Observa -se que nesses intervalos estão computadas as ocorrências de respostas em branco para as três questões discursivas de Conhecimento Específico.
A análise de cada uma destas questões será feita a seguir. 6.3.2.1 Análise da Questão Discursiva 3 do Componente de Conhecimento Específico A questão discursiva 3 de Conhecimentos específicos foi a mesma para estudantes do Bacharelado e da Licenciatura.
Nessa questão, cujos resultados aferidos encontram -se descritos nas Tabelas 6.18a e 6.18b, respectivamente para os cursos de Licenciatura e de Bacharelado, as Médias dos estudantes de cursos de Licenciatura foram 2,6, 0,9, respectivamente, para os cursos de Educação Presencial e os a Distância ( 2,1 para o total ), ao passo que a Média dos estudantes de Bacharelado foi 24,6.
Esta foi a questão de mais baixo desempenho dentre as três discursivas de Conhecimentos Específicos, para as duas Modalidades da Licenciatura e para o Bacharelado.
Para todas as combinações de Habilitação e Modalidade, a nota Mínima foi 0,0, as mesmas em todas as regiões do Brasil.
A nota Máxima foi 100,0 para a Educação Presencial da Licenciatura e para o Bacharelado.
Na Educação a Distância da Licenciatura, a nota Máxima foi 60,0.
O Gráfico 6.11 mostra a distribuição das notas na questão discursiva 3, do Componente de Conhecimento Específico da área de Ciência da Computação ( Licenciatura/Bacharelado ).
O intervalo modal, para ambas as Habilitações é aquele formado pelos participantes que deixaram a resposta à questão em branco, com mais de 60 \% na Licenciatura e cerca de 30 \% no Bacharelado. 6.3.2.2 Comentários sobre as respostas à Questão Discursiva 3 A Questão 3, proposta para os estudantes das duas habilitações, abordava mais de um assunto importante para o profissional de Computação : capacidade de ler e interpretar um algoritmo recursivo, estruturas de dados, análise de complexidade e linguagens de programação.
São assuntos que constam das Diretrizes Curriculares Nacionais, do Currículo de Referência da SBC ( Sociedade Brasileira de Computação ) e dos programas dos principais cursos de graduação em computação.
O padrão de resposta cobria possibilidades de utilização de diversas linguagens de programação para as respostas dos estudantes.
A questão não admitia mais de uma interpretação e estava redigida em linguagem clara e acessível ao nível dos estudantes.
Dessa forma, as respostas dos estudantes mantiveram -se em conformidade com as opções previstas no padrão de resposta.
Pode ser considerada uma questão de dificuldade média, principalmente porque exigia conhecimento de mais de um assunto para resolvê -la.
O item ‘ a ’ tratava da complexidade do algoritmo dado e obteve quantidade de respostas e acertos dentro do esperado.
Grande parte dos erros refletiu dificuldade com a notação e a compreensão do algoritmo, mesmo em alunos que ganharam pontuação integral, foi difusa e o que denota uma dificuldade grande no tópico.
O item ‘ b ’ demonstrou outra faceta - a do desconhecimento de estruturas de dados e de como tratar recursão.
As duas estruturas sobre as quais o problema exigia tratamentos diferentes foram manipuladas sem distinção, misturando conceitos e revelando falta de familiaridade com os assuntos.
O vocabulário relativo aos conteúdos é conhecido, mas houve dificuldades com a notação e a sintaxe das linguagens e, muitas vezes, raciocínios aparentemente corretos eram expressos de maneira equivocada.
Em torno de 65 \% dos estudantes de Licenciatura não tentaram responder à questão.
Dentre os que o fizeram, as respostas eram basicamente incorretas em todos os aspectos, demonstrando deficiência muito grande de formação nos assuntos avaliados.
Não houve procedimentos mais frequentes a destacar – os acertos foram poucos e os erros variados e generalizados.
Já no Bacharelado, embora com resultado muito aquém do desejado, houve uma concentração de erros em dois assuntos : no entendimento da recursão e na manipulação da lista.
As abordagens dos estudantes que tentaram resolver a questão seguiram o previsto no padrão de resposta e, com grande regularidade, respondiam ao item ‘ a ’.
No item ‘ b ’ faziam - errando ou não - a varredura da lista, as duas operações previstas para realizar com a pilha : empilhar e desempilhar, e a impressão.
Um número grande de alunos simplesmente replicou o algoritmo original ou então usou soluções ( corretas ou não ) que continham algum tipo de recursão, o que não era permitido pelo enunciado.
Outro conjunto grande de alunos tratou a fila como se fosse uma pilha, o que é um erro conceitual grave.
Houve também uma demonstração clara da falta de entendimento do conceito de ponteiro.
As respostas que receberam as notas mais baixas continham erro no item sobre complexidade ( item ‘ a ’ ) e também erros sérios ao se escrever o algoritmo solicitado, como tratar a lista como pilha ( item ‘ b ’ ).
Nas respostas medianas, normalmente, o item ‘ a ’ foi respondido corretamente, mas havia erro sério no item ‘ b ’, normalmente no primeiro laço previsto nos padrões de resposta, a varredura da lista.
Nas respostas boas, normalmente o item ‘ a ’ estava correto, assim como também grande parte do item ‘ b ’ ; varrendo a lista e manipulando a pilha, mas errando a manipulação de ponteiros, deixando de imprimir ou não invertendo a impressão.
A deficiência nas áreas de programação, recursão, estruturas de dados e análise de complexidade, assuntos que foram avaliados pela Questão 3, mostrou -se enorme.
A resposta dos participantes do Enade revelou falta de conhecimentos essenciais, como manipulação de ponteiros e memória, talvez porque linguagens atuais como Java e Python omitam esses conceitos.
Isso pode não representar um impedimento na realização das atividades cotidianas, mas gera uma ofuscação dos conceitos, através da introdução de uma camada de abstração que impede que eles desenvolvam tipos mais avançados de aplicativos, como drivers ou sistemas operacionais. 6.3.2.3 Análise da Questão Discursiva 4 do Componente de Conhecimento Específico As Tabelas 6.19a e 6.19b contêm as informações relativas à questão 4 do conjunto de questões do Componente de Conhecimento Específico.
A questão 4 não foi a mesma para as duas Habilitações.
Para ambas as Habilitações, o desempenho dos estudantes nessa questão foi superior ao da questão 3 ( a mesma para todos os estudantes ) e inferior ao das questões 5 ( também diferente para cada Habilitação ).
As Médias para os alunos de Licenciatura ( ver Tabela 6.19a ) foram, respectivamente, para os de cursos de Educação Presencial e de Educação a Distância, 8,3 e 6,6 ( 7,8 para o total ).
A nota Máxima ( 95,0 ) da Educação Presencial da Licenciatura foi atingida por, pelo menos, um concluinte da região Nordeste.
Na Educação a Distância, a nota Máxima foi 80,0, na região Sudeste.
A Mediana foi zero para ambas as Modalidades de ensino da Licenciatura em todas as regiões.
A nota Mínima foi zero para todas as combinações de Modalidade de ensino e região.
A Média para os alunos de Bacharelado ( ver Tabela 6.19b ) foi 26,0.
A nota Máxima ( 100,0 ) foi atingida por, pelo menos, um concluinte de todas as regiões, com exceção da região Centro-Oeste, onde a nota Máxima foi 95,0.
A Mediana foi zero para todas as regiões, com exceção da região Sul ( 20,0 ), indicando que mais do que 50 \% dos estudantes tiraram 0,0.
A nota Mínima foi zero para todas as regiões.
O Gráfico 6.12 representa a distribuição de notas da questão discursiva 4, do Componente de Conhecimento Específico.
Essa distribuição tem moda na classe de estudantes que deixaram a resposta em branco, tanto na Habilitação Bacharelado quanto na Licenciatura. 6.3.2.4 Comentários sobre as respostas à Questão Discursiva 4 O tema da questão é de importância fundamental nos cursos de computação.
A Questão 4 tratava da modelagem de sistemas, o que faz parte das Diretrizes Curriculares Nacionais, do Currículo de Referência da SBC ( Sociedade Brasileira de Computação ) e dos programas dos principais cursos de graduação em computação.
A questão pode ser considerada de dificuldade média para fácil, embora exigisse o conhecimento de notação específica.
O enunciado estava bastante claro e com informações e dados suficientes para a solução correta.
O conhecimento exigido para a solução não ultrapassava o conhecimento básico da área.
Ficou claro que todos sabiam o que foi \textbf{indagado}, apesar de não conseguirem responder corretamente ao que foi solicitado.
O padrão de resposta cobriu as respostas possíveis com grande abrangência ; não houve quantidade significativa de acertos com estrutura não prevista no padrão.
As respostas, inclusive as erradas, seguiam o padrão de resposta, mesmo quando usavam linguagem diferente da UML.
As respostas dos estudantes do curso de Bacharelado corresponderam ao previsto no padrão.
No entanto, principalmente na Licenciatura, além do número excessivo de respostas em branco, muitos evidenciaram desconhecimento do assunto, além do esperado.
As respostas insuficientes foram muito pouco além da identificação das classes e de seus relacionamentos.
As respostas dos estudantes de Licenciatura, quando ocorreram, foram dadas de forma bastante insatisfatória, sem sequer ser utilizada a notação correta da linguagem UML, solicitada no enunciado.
A UML, linguagem indicada para resolver o problema, deveria ser conhecida e utilizada pelos alunos.
No entanto, ocorreram muitos erros, principalmente na Licenciatura.
Algumas respostas erradas mantiveram um padrão de erro, usando outra linguagem, criando classes inexistentes, não identificando corretamente herança.
O conjunto de respostas da Licenciatura e do Bacharelado demonstrou uma diferença grande nos resultados.
Não houve qualquer semelhança positiva no tratamento da questão entre as respostas dos estudantes dos dois cursos.
Os estudantes de Licenciatura demonstraram desconhecimento quase que absoluto do assunto.
A abordagem básica, tanto no Bacharelado, quanto dentre aqueles que responderam à questão de forma aceitável no curso de Licenciatura ( cerca de 5 \% ), foi a abordagem comumente adotada nesse tipo de questão : identificação das classes e relacionamentos para, em seguida, anotar atributos, métodos e heranças.
Nas melhores respostas os estudantes encontraram as classes envolvidas e os relacionamentos básicos.
Além disso, também identificaram os atributos.
Os problemas sérios ocorreram na identificação da herança e, principalmente, dos métodos.
Os erros mais comuns no Bacharelado foram : i ) atribuir à classe"Departamento"os métodos que deveriam ter sido designados às subclasses da classe"Colaborador"; ii ) dificuldades na identificação e representação do conceito de herança ; iii ) identificação de outros elementos no modelo ( tais como atributos ou métodos equivocados na classe"Departamento").
No curso de Licenciatura, ficou demonstrado que os formandos, salvo raras exceções, têm apenas rudimentos de modelagem conceitual.
Representaram a classe"Colaborador"e, eventualmente, associaram essa classe a"Departamento".
As respostas no Bacharelado mantiveram uma estrutura básica e, dentre os que responderam à questão, houve soluções bem estruturadas, com notas próximas ou iguais à nota máxima.
Na Licenciatura, também não houve diversidade, mas as respostas eram quase ou totalmente erradas.
As respostas com as notas mais baixas continham erros inclusive na identificação das classes, problema considerado, no caso dessa questão, bastante simples.
As respostas medianas foram aquelas com acerto na identificação das classes, dos relacionamentos entre elas e dos atributos.
Duas deficiências foram mais fortemente reveladas.
A mais relevante foi a falta de capacidade de modelar um problema simples, e a segunda, o desconhecimento da linguagem indicada para a solução ( UML ), bastante importante atualmente.
Esse desconhecimento foi muito mais evidente na Licenciatura, cujos estudantes não identificavam sequer a totalidade das classes, quando não indicavam classes inexistentes.
Os estudantes de Bacharelado, mesmo errando, identificavam as classes e usavam a UML. 6.3.2.5 Análise da Questão Discursiva 5 do Componente de Conhecimento Específico As Tabelas 6.20a e 6.20b contêm, respectivamente para Licenciatura e Bacharelado, as informações relativas às questões 5 do Componente de Conhecimento Específico.
A questão 5 do Bacharelado e da Licenciatura eram distintas.
O desempenho dos estudantes de ambas as Habilitações nessas questões foi o mais alto dentre as três questões discursivas desse componente.
As Médias para os alunos de Licenciatura ( ver Tabela 6.20a ) foram, respectivamente, para os de cursos de Educação Presencial e de Educação a Distância, 40,4 e 36,5 ( 39,1 para o total ).
A Mediana foi 42,5 para a Educação Presencial, e 30,0 para a Educação a Distância.
A nota Máxima ( 100,0 ) foi atingida em todas as regiões por, pelo menos, um participante de curso de Educação Presencial de Licenciatura.
Também o Nordeste e o Sudeste obtiveram nota Máxima 100,0 na Educação a Distância da Licenciatura.
A nota Mínima foi zero para todas as combinações de Modalidade de ensino e região, com exceção da região Sul, onde a nota Mínima foi 50,0 para cursos a Distância.
A Média para os alunos de Bacharelado ( ver Tabela 6.20b ) foi 34,7.
A nota Máxima ( 100,0 ) foi atingida por, pelo menos, um concluinte de todas as regiões.
A Mediana foi 20,0 para o Brasil e para as regiões Nordeste e Sudeste.
A Mediana nas regiões Norte e Centro-Oeste foi zero, e na região Sul foi 30,0.
A nota Mínima foi zero para todas as regiões.
O Gráfico 6.13 representa a distribuição das notas da questão discursiva 5 do Componente de Conhecimento Específico para o conjunto de estudantes do Bacharelado e da Licenciatura.
Essa distribuição tem moda no intervalo correspondente ao grupo que deixou a resposta em branco, com frequência em torno de 30 \% para ambas as Habilitações. 6.3.2.6 Comentários sobre as respostas à Questão Discursiva 5 Questão 5 – Bacharelado A Questão 5, específica do Bacharelado, tratava de um algoritmo de busca em grafos, assunto relevante para várias disciplinas da área da computação.
Os conteúdos avaliados pela questão fazem parte das Diretrizes Curriculares Nacionais para a área, do currículo referência da SBC ( Sociedade Brasileira de Computação ) e da grande maioria dos bons cursos de graduação em computação.
O enunciado estava bastante claro e com informações e dados suficientes para a solução correta.
O conhecimento exigido para a solução não ultrapassava o conhecimento básico da área.
O item ‘ a ’ dessa questão pode ser considerado de dificuldade média para fácil, embora exigisse o conhecimento de um algoritmo específico.
O item ‘ b ’ pode ser considerado muito fácil, mas surpreendentemente foi respondido com menos acertos do que o item ‘ a ’.
O padrão de resposta cobria as respostas possíveis com grande abrangência e proporcionava boa discriminação ; não houve quantidade significativa de acertos com estrutura não prevista no padrão.
As respostas, inclusive as erradas, buscavam seguir o procedimento do padrão de resposta adotado.
As respostas corretas do item ‘ a ’ seguiram maciçamente a primeira das soluções do padrão, a mais óbvia.
As respostas erradas foram consequência do emprego equivocado do algoritmo.
No item ‘ b ’, as respostas ou seguiram o padrão e acertaram ou divergiram do padrão e erraram.
Nesse item, surpreendentemente, os erros evidenciaram o desconhecimento do que seja uma matriz de adjacência, assunto tratado em várias disciplinas.
Houve duas formas mais comuns de erro, ora marcando ligações inexistentes, ora deixando de marcar as existentes.
A linguagem utilizada foi adequada porque era muito restrita.
O item ‘ a ’ foi majoritariamente respondido seguindo as setas em sequência alfabética, com alguns sabendo o que fazer quando chegava a um loop, e outros, não.
O erro principal foi não saber como voltar atrás quando um loop era encontrado.
Para o item ‘ b ’, o mais frequente foi o preenchimento da matriz com uns e zeros.
Nesse item houve muitas respostas em branco, apesar de ser considerado um assunto de solução fácil.
Foi surpreendente a quantidade de brancos e erros.
As soluções corretas variaram pouco.
Das quatro possibilidades, quase sempre a adotada foi a primeira explicitada no padrão de respostas.
Com as respostas erradas, houve alguma variação, a maioria sendo uma solução parcial que parava no loop.
Os alunos de Bacharelado tiveram um desempenho aquém do desejado, embora um número expressivo tenha tirado nota máxima.
Erraram na busca ( item ‘ a ’ ), ao pararem em determinado nó, por não saberem como continuar e/ou por misturarem os critérios usados para continuar ao encontrarem um loop ; ora usavam um critério, ora outro.
Além disso, o erro frequente no item ‘ b ’, totalmente inesperado, reduziu a quantidade de notas maiores do que 70.
As respostas com as notas mais baixas estavam corretas em um trecho da busca ou na matriz.
Nas respostas medianas, a busca, ou trechos da busca, e a matriz estavam corretos.
As respostas boas acertaram a busca, ou grande parte dela, e a matriz.
Embora significativa parcela dos estudantes tenha acertado a questão integralmente, a quantidade de brancos e de notas fracas revelou o desconhecimento do conceito bastante elementar de matriz de adjacência.
Além disso, grande parte das respostas ao item ‘ a ’ iam até o ponto da busca em que era necessário tomar uma decisão.
A parada nesse ponto demonstra desconhecimento da parte principal do algoritmo.
Essas falhas são relevantes, dada a importância do assunto no currículo do curso.
Questão 5 – Licenciatura A questão avaliou o conhecimento sobre aspectos \textbf{pedagógicos} de três tipos de software educativo, temática bastante relevante para a habilitação em Licenciatura de Ciência da Computação.
O assunto consta das Diretrizes Curriculares Nacionais para a área e da grande maioria dos bons cursos de Licenciatura da Área.
Era uma questão simples, redigida com clareza, linguagem e nível de profundidade adequados.
O enunciado fornecia informações e dados suficientes para a solução correta da questão.
Dessa forma, não havia possibilidade de interpretação diferente da esperada.
O comando solicitava identificar aspectos \textbf{pedagógicos} de dois tipos de software dos três citados no enunciado.
A resposta deveria explicitar dois aspectos que justificassem a escolha e um aspecto limitante da eficácia do software escolhido.
A questão pode ser considerada de dificuldade média para fácil.
O que a questão solicitava era básico para a área de ensino em computação, embora exigisse o conhecimento das características de três tipos de software.
As respostas dos estudantes, inclusive as erradas, seguiram o padrão esperado e mantiveram -se, quase sempre, num \textbf{subconjunto} do universo previsto, no entanto, com o uso de vocabulário menos sofisticado.
As respostas não previstas no padrão, e aceitas, beiravam o senso comum, espaço aberto pela amplitude da pergunta.
As respostas corretas, muitas vezes, continham exemplo de algum tipo de software ( simulador, micromundo ou jogo ), o que não foi pedido na questão.
As incompletas citavam um tipo de software somente ou deixavam de elencar as deficiências.
Os tipos de software mais frequentemente citados foram os jogos e os simuladores.
Os respondentes valorizaram a característica lúdica dos jogos e a possibilidade de tratar de problemas do mundo real com procedimentos em simuladores.
Um erro comum foi associar o software educativo aberto ao software livre, reforçando a resposta com exemplo, descrição e potencial \textbf{pedagógico} de algum software desse tipo.
Muitas vezes, os softwares eram descritos, mas as propostas \textbf{pedagógicas} não eram citadas.
Em algumas respostas, havia menção a dois softwares distintos, mas sem diferenciação do limite e potencial de cada um, anunciando as mesmas características para ambos.
Houve também dissertação sobre construtivismo e sua relação com a educação, aprendizagem e papel docente, temas que não faziam parte da pergunta.
A limitação dos softwares, muitas vezes, recaiu sobre as limitações da infraestrutura das escolas e não naquelas do próprio software.
Repetiram -se respostas que, embora se restringissem ao senso comum, atendiam fracamente ao que foi solicitado pelo comando da questão.
O resultado ficou muito abaixo do esperado de formandos de um curso de Licenciatura em computação, principalmente se considerarmos a simplicidade da questão.
Aparentemente, a pergunta não foi lida com o devido cuidado, e muitos estudantes apresentaram apenas uma vantagem para cada tipo de software ; outros não citaram desvantagens e muitos não escolheram dois tipos de software, e sim, dois softwares de um mesmo tipo.
Além disso, houve mistura entre os conceitos de software aberto e software livre, inadmissível para licenciandos de Ciência da Computação.
As respostas com as notas mais baixas, quase sempre, citavam apenas um tipo de software e não relatavam, principalmente, as deficiências.
As respostas medianas, normalmente, citavam os dois tipos de software e seus aspectos \textbf{pedagógicos} positivos.
As respostas boas eram bem organizadas e citavam quase tudo o que era pedido.
A deficiência mais marcante foi, sem dúvida, a dificuldade com a linguagem da Pedagogia - os pontos positivos ou negativos enunciados, na maioria das vezes, foram elencados com o uso de vocabulário não específico. 6.3.3 Considerações Finais O resultado geral ficou aquém do desejado, principalmente no curso de Licenciatura.
As dificuldades dos estudantes deste curso ficaram patentes nos resultados das questões 3 e 4.
Nas duas habilitações, tais questões trataram de assuntos essenciais para a construção de qualquer software, e os respondentes demonstraram lacunas sérias de formação, a destacar a dificuldade no trato da recursão e da modelagem de um sistema.
A questão 5 do Bacharelado acrescentou, a essa lista, o desconhecimento dos detalhes da aplicação do algoritmo de busca e, surpreendentemente, o da matriz de adjacência.
A questão sobre software educativo - exclusiva da Licenciatura - demonstrou falta de conhecimento mais sofisticado do vocabulário da área de ensino e, além disso, aspectos \textbf{pedagógicos} foram substituídos pelo senso comum.

% matched lemmas: indagar, inverso, pedagógico, subconjunto
\end{textsample}
